% =========================================================================
% Problem Formulation + Proposed Architecture (extend_GNN -> "DockGNN")
% -------------------------------------------------------------------------
% Usage: % =========================================================================
% Problem Formulation + Proposed Architecture (extend_GNN -> "DockGNN")
% -------------------------------------------------------------------------
% Usage: % =========================================================================
% Problem Formulation + Proposed Architecture (extend_GNN -> "DockGNN")
% -------------------------------------------------------------------------
% Usage: % =========================================================================
% Problem Formulation + Proposed Architecture (extend_GNN -> "DockGNN")
% -------------------------------------------------------------------------
% Usage: \input{formulation_and_architecture} inside an IEEEtran document.
% Requires in the preamble:
%   \usepackage{amsmath,amssymb,bm}
%   \usepackage{booktabs,array,multirow,graphicx}
%   \usepackage{algorithm}
%   \usepackage{algpseudocode}
% Citation keys are placeholders -- suggested BibTeX entries are listed in
% comments at the end of this file. Replace \cite{...} keys with your own.
% "DockGNN" is a placeholder model name; rename globally if you prefer.
% Contains four sections: Problem Formulation, Architecture, Surrogate
% Calibration (sec:calibration), and Policy Optimization (sec:training).
% =========================================================================

\section{Problem Formulation}
\label{sec:problem}

\begin{table*}[t]
\caption{Summary of Notation}
\label{tab:notation}
\centering
\footnotesize
\begin{tabular}{@{}l p{0.295\textwidth} @{\qquad} l p{0.295\textwidth}@{}}
\toprule
Symbol & Meaning & Symbol & Meaning \\
\midrule
\multicolumn{2}{@{}l}{\itshape Problem setting} &
\multicolumn{2}{@{}l}{\itshape Architecture} \\
$\mathcal{W}$ & grid workspace (four-connected unit cells) &
  $\bm{x}^{A}_k,\bm{x}^{J}_j,\bm{x}^{D}_d$ & AMR / job / dock input features ($\mathbb{R}^{12},\mathbb{R}^{11},\mathbb{R}^{9}$) \\
$\Delta(u,v)$ & Manhattan distance between positions $u,v$ &
  $\bm{x}^{\mathrm{act}}_{\omega,k,j}$ & per-action decision features ($\mathbb{R}^{3}$) \\
$\mathcal{J},\ n$ & job set; number of jobs &
  $h$ & hidden width of all embeddings ($h{=}128$) \\
$\mathcal{M},\ m$ & AMR fleet; fleet size &
  $\bm{h}^{A}_k,\ \bm{h}^{D}_d$ & AMR and dock embeddings \\
$\mathcal{T}$ & set of material types &
  $\bar{\bm{h}}^{J}_j$ & job embedding before dock fusion \\
$\mathcal{D}^{\mathrm{in}},\mathcal{D}^{\mathrm{out}},\mathcal{D}$ & inbound docks; outbound docks; their union &
  $\bm{H}^{A},\bm{H}^{D}$ & stacked AMR / dock embedding matrices \\
$\tau_j$ & material type of job $j$ &
  $\mathrm{MHA},\mathrm{LN}$ & multi-head attention; layer normalization \\
$o_j,\ d_j$ & inbound and outbound dock of job $j$ &
  $\bm{g}^{(l)}_j$ & job embedding after fusion and $l$ GIN layers \\
$r_j$ & release time of job $j$ &
  $\mathcal{G}_t=(\mathcal{J},E_t)$ & contention--precedence graph at step $t$ \\
$p_j=p(\tau_j)$ & dock service duration of job $j$ &
  $\mathcal{N}(j),\ \deg_j$ & graph neighbors and degree of job $j$ \\
$\pi_j,\ \delta_j$ & pickup and delivery operations of job $j$ &
  $\epsilon^{(l)},\ L$ & learnable GIN self-weight; number of GIN layers \\
$Q$ & per-type AMR carrying capacity &
  $\bm{e}_{\omega}$ & learned operation-type embedding \\
$\bm{a},\ a_j$ & AMR assignment vector; AMR assigned to job $j$ &
  $d(\omega,j)$ & dock served by the action ($o_j$ if pickup, else $d_j$) \\
$\sigma,\ \sigma_t$ & operation sequence; prefix committed up to step $t$ &
  $z_{\omega,k,j}$ & logit of action $(\omega,k,j)$ \\
$s^{\pi}_j,c^{\pi}_j$ & start / completion of pickup service of job $j$ &
  $V_{\phi}$ & critic (state-value) network, parameters $\phi$ \\
$s^{\delta}_j,c^{\delta}_j$ & start / completion of delivery service of job $j$ &
  & \\
$t_k(u,v)$ & realized travel time of AMR $k$ from $u$ to $v$ &
  \multicolumn{2}{@{}l}{\itshape Training} \\
$\Phi$ & collision-aware executor &
  $R,\ R(\sigma_{1:t})$ & dispatching rule; completion of prefix $\sigma_{1:t}$ by $R$ \\
$C_{\max},\ F_k$ & makespan; finish time of AMR $k$ incl.\ depot return &
  $C[\cdot]$ & shorthand for $C_{\max}(\Phi(\cdot))$ \\
$W_d$ & number of physical waiting cells at dock $d$ &
  $A_t,\ \hat{A}_t$ & stepwise counterfactual advantage; standardized \\
\multicolumn{2}{@{}l}{\itshape Sequential decision process} &
  $A^{\mathrm{lb}}_t$ & load-balance shaping advantage \\
$s_t,\ a_t$ & state and action at decision epoch $t$ &
  $\lambda_{\mathrm{lb}},\ \lambda_{H}$ & shaping and entropy coefficients \\
$\omega$ & operation type, $\omega\in\{\textsc{pick},\textsc{deliver}\}$ &
  $\mathcal{H}_t$ & entropy of the masked policy at step $t$ \\
$\Lambda(s_t)$ & feasible-action set (mask) in state $s_t$ &
  $\mathcal{L}(\theta)$ & training loss \\
$\hat{\Psi}$ & calibrated surrogate transition model &
  $B,\ E$ & batch size (episodes); number of training epochs \\
$\mathcal{E}_d$ & committed service events $(s_e,c_e,j_e,\omega_e)$ at dock $d$ &
  $T_{\mathrm{val}},\ \kappa$ & validation interval; infeasibility penalty \\
$\hat{b}_k,\ x_k$ & next-available time and position of AMR $k$ &
  $\mathbf{1}[\cdot]$ & indicator function \\
$\rho$ & ready time of an AMR at a dock &
  & \\
$\hat{t}(u,v)$ & calibrated travel-time estimate \eqref{eq:travelcal} &
  & \\
$\hat{w}_d(\rho)$ & queue-aware dock-wait estimate \eqref{eq:waitcal} &
  & \\
$q_d(\rho)$ & services still unfinished at time $\rho$ at dock $d$ &
  & \\
$\alpha,\ \beta$ & travel-calibration scale and offset &
  & \\
$\varphi(\cdot)$ & dock-wait penalty table over queue depths &
  & \\
$\pi_{\theta}$ & scheduling policy with parameters $\theta$ &
  & \\
\bottomrule
\end{tabular}
\end{table*}

\subsection{System Description}

We consider the internal transfer problem of a cross-docking terminal
served by a fleet of autonomous mobile robots (AMRs). Table~\ref{tab:notation}
summarizes the notation used throughout the paper. The terminal floor
is modeled as a four-connected unit grid $\mathcal{W}$ in which a set of
\emph{inbound docks} $\mathcal{D}^{\mathrm{in}}$ is located along one
boundary and a set of \emph{outbound docks} (processing stations)
$\mathcal{D}^{\mathrm{out}}$ along the opposite boundary; the AMR depots
lie between them. Consistent with cross-docking practice, the floor is
free of static obstacles, so congestion arises from vehicle interference
and dock contention rather than from the layout itself. AMRs travel at
unit speed between adjacent cells, and the distance between two
positions $u,v$ is the Manhattan metric $\Delta(u,v)$.

A set of transfer jobs $\mathcal{J}=\{1,\dots,n\}$ must be moved across
the terminal by a fleet $\mathcal{M}=\{1,\dots,m\}$ of homogeneous AMRs.
Each job $j\in\mathcal{J}$ represents one unit load of a material type
$\tau_j\in\mathcal{T}$ that becomes available at its inbound dock
$o_j\in\mathcal{D}^{\mathrm{in}}$ at release time $r_j\ge 0$ and must be
delivered to a designated outbound dock
$d_j\in\mathcal{D}^{\mathrm{out}}$. Handling a unit load of type $\tau$
occupies a dock for a type-dependent service duration $p(\tau)$; we
write $p_j=p(\tau_j)$ and assume the loading service at the inbound dock
and the unloading/processing service at the outbound dock take the same
time, reflecting symmetric handling of a unit load. Each job therefore
consists of an ordered pair of operations: a \emph{pickup} operation
$\pi_j$ executed at $o_j$ and a \emph{delivery} operation $\delta_j$
executed at $d_j$, both performed by the same AMR (no transshipment).
An AMR may consolidate loads subject to a per-type carrying capacity
$Q$: at no time may it carry more than $Q$ units of any single material
type.

Docks are exclusive single-server resources: at most one AMR may be in
service at a dock at any time, and services are non-preemptive. An AMR
arriving at an occupied dock queues in a bounded set of waiting cells
adjacent to that dock; when the waiting line is saturated, the AMR must
hold at an upstream position and re-approach later. Finally, AMR motion
must be collision-free: at any unit time step, each grid cell may be
occupied by at most one AMR.

\subsection{Formal Statement}

A solution is a pair $(\bm{a},\sigma)$, where
$\bm{a}=(a_1,\dots,a_n)\in\mathcal{M}^{n}$ assigns each job to an AMR
and $\sigma$ is a total order over the $2n$ operations
$\{\pi_j,\delta_j\}_{j\in\mathcal{J}}$ that determines the sequence in
which operations are committed. Let $s^{\pi}_j,c^{\pi}_j$
(resp.\ $s^{\delta}_j,c^{\delta}_j$) denote the service start and
completion times of $\pi_j$ (resp.\ $\delta_j$) induced by executing
$(\bm{a},\sigma)$. A feasible execution satisfies:
\begin{align}
c^{\pi}_j &= s^{\pi}_j + p_j, \qquad
c^{\delta}_j = s^{\delta}_j + p_j,
&& \forall j\in\mathcal{J}, \label{eq:nonpreempt}\\
s^{\pi}_j &\ge r_j,
&& \forall j\in\mathcal{J}, \label{eq:release}\\
s^{\delta}_j &\ge c^{\pi}_j,
&& \forall j\in\mathcal{J}, \label{eq:precedence}
\end{align}
\begin{align}
\big|\{\,j : a_j{=}k,\ \tau_j{=}\tau,\ c^{\pi}_j \le u < c^{\delta}_j \}\big| &\le Q,
&& \forall k,\tau,u, \label{eq:capacity}\\
\big[s_e, c_e\big) \cap \big[s_{e'}, c_{e'}\big) &= \emptyset,
&& \forall e\ne e' \ \text{at the same dock}, \label{eq:exclusive}
\end{align}
where \eqref{eq:nonpreempt} enforces non-preemptive services,
\eqref{eq:release} release dates, \eqref{eq:precedence} pickup-before-%
delivery precedence, \eqref{eq:capacity} the per-type carrying capacity,
and \eqref{eq:exclusive} dock exclusivity over all service events $e$.
In addition, for any two operations $u\prec v$ executed consecutively by
the same AMR $k$,
\begin{equation}
s_v \;\ge\; c_u + t_k\big(\mathrm{loc}(u),\mathrm{loc}(v)\big),
\label{eq:travel}
\end{equation}
where $t_k(\cdot,\cdot)\ge\Delta(\cdot,\cdot)$ is the realized travel
time of AMR $k$, which may exceed the Manhattan distance due to
collision avoidance, queuing detours, and dock-clearing maneuvers.

The realized travel and waiting times are not fixed constants but are
determined by a high-fidelity \emph{executor} $\Phi$, which routes every
AMR leg with space--time $\mathrm{A}^{*}$ search over a shared
reservation table (at most one AMR per cell per time step), manages the
bounded dock waiting lines, and inserts upstream holding maneuvers when
lines are saturated. The problem is therefore a simulation-based
optimization: find
\begin{equation}
(\bm{a}^{*},\sigma^{*}) \;=\;
\arg\min_{(\bm{a},\sigma)} \; C_{\max}\big(\Phi(\bm{a},\sigma)\big),
\label{eq:objective}
\end{equation}
where the makespan $C_{\max}=\max_{k\in\mathcal{M}} F_k$ and $F_k$ is
the time at which AMR $k$ completes its final activity, including its
return trip to the depot. The problem generalizes both the capacitated
pickup-and-delivery problem and parallel-machine scheduling with
sequence-dependent setups, and is therefore NP-hard; the coupling to
collision-free multi-agent routing further compounds its difficulty
\cite{stern2019mapf}.

\subsection{Sequential Decision Formulation}
\label{sec:mdp}

We cast the construction of $(\bm{a},\sigma)$ as an episodic sequential
decision process with exactly $2n$ decision epochs, one per committed
operation. At epoch $t$, the state
\begin{equation}
s_t=\big(\{\bm{x}^{A}_{k}\},\, \{\bm{x}^{J}_{j}\},\, \{\bm{x}^{D}_{d}\},\,
\{\mathcal{E}_d\},\, \sigma_{t}\big)
\label{eq:state}
\end{equation}
collects the AMR states (position, next-available time, per-type
inventory), the job states (unpicked / onboard / delivered, carrier
identity), the dock states, the set of committed dock service events
$\mathcal{E}_d=\{(s_e,c_e,j_e,\omega_e)\}$ at each dock $d$, and the
partial operation sequence $\sigma_t$.

An action is a triple
\begin{equation}
a_t=(\omega,k,j)\in\{\textsc{pick},\textsc{deliver}\}\times\mathcal{M}\times\mathcal{J},
\label{eq:action}
\end{equation}
giving a flat action space of size $2mn$ restricted by a feasibility
mask $\Lambda(s_t)$: $(\textsc{pick},k,j)$ is legal iff $j$ is unpicked
and AMR $k$ carries fewer than $Q$ units of type $\tau_j$;
$(\textsc{deliver},k,j)$ is legal iff $j$ is onboard AMR $k$. The mask
guarantees by construction that every complete episode yields a feasible
solution, so no repair mechanism is required.

Transitions follow a deterministic surrogate model $\hat{\Psi}$ that
advances the acting AMR's clock using a \emph{calibrated} travel-time
estimate $\hat{t}(\cdot,\cdot)$ and a queue-aware dock-wait estimate
$\hat{w}_d(\cdot)$ fitted offline against the executor $\Phi$
(Section~\ref{sec:calibration}); committing an action appends the
corresponding service window to $\mathcal{E}_d$. For a pickup, the
service start of job $j$ by AMR $k$ located at $x_k$ with next-available
time $\hat{b}_k$ is
\begin{equation}
s^{\pi}_j=\underbrace{\max\!\big(\hat{b}_k+\hat{t}(x_k,o_j),\, r_j\big)}_{\text{ready time }\rho}
\;+\;\hat{w}_{o_j}(\rho),
\label{eq:surrogate}
\end{equation}
and analogously for deliveries at $d_j$. After $2n$ steps the completed
solution is evaluated by the executor, and the learning objective is to
minimize the expected executed makespan
$\mathbb{E}_{\pi_\theta}\!\left[C_{\max}(\Phi(\bm{a},\sigma))\right]$;
the policy-gradient training procedure and its dispatch-rule baseline
are described in Section~\ref{sec:training}.

% =========================================================================

\section{Dock-Aware Graph Scheduler}
\label{sec:architecture}

\subsection{Overview and Design Rationale}

Existing learning-to-dispatch architectures for shop scheduling encode
jobs (operations) as graph nodes and machines as flat feature vectors
\cite{zhang2020l2d,song2023fjsp}. Applied to cross-docking, this design
has three blind spots. First, the binding resources are the docks: with
$m$ AMRs contending for $|\mathcal{D}^{\mathrm{in}}|+|\mathcal{D}^{\mathrm{out}}|$
single-server docks, queuing at docks---not processing---dominates the
makespan, yet job-only graphs represent dock congestion at best as
scalar features replicated across jobs. Second, the temporal structure
of dock occupancy (which service windows are already committed, and
when) is invisible. Third, the AMR fleet is encoded independently
per vehicle, so mutual interference cannot be anticipated.

The proposed scheduler, hereafter \textsc{DockGNN}, addresses all three
by treating docks as first-class graph entities. It comprises
(i)~heterogeneous encoders for AMR, job, and dock nodes;
(ii)~an AMR--dock interaction block that exchanges information between
the fleet and the dock set;
(iii)~a job \emph{contention--precedence} graph processed by
degree-normalized graph isomorphism network (GIN) layers
\cite{xu2019gin}; and
(iv)~a factored actor head that scores each masked action from the
operation type, the acting AMR, the target job, the operation-specific
dock, and a small set of per-action decision features. The full model
has $6.2\times10^{5}$ parameters ($5.7\times10^{5}$ actor,
$4.9\times10^{4}$ critic) with hidden width $h=128$ throughout.

% TODO: architecture overview figure
% \begin{figure*}[t]
%   \centering
%   \includegraphics[width=\textwidth]{figures/dockgnn_architecture}
%   \caption{Overview of \textsc{DockGNN}. Dock, AMR, and job features are
%   encoded separately; an AMR--dock attention block exchanges congestion
%   information; jobs fuse their inbound/outbound dock embeddings and pass
%   messages over the contention--precedence graph; the actor scores each
%   feasible (operation, AMR, job) triple.}
%   \label{fig:architecture}
% \end{figure*}

\subsection{Heterogeneous State Representation}
\label{sec:features}

At every decision epoch the state \eqref{eq:state} is rendered into four
feature families, summarized in Table~\ref{tab:features}. All features
are normalized to comparable magnitudes by fixed scales.

\emph{Dock nodes.} Each dock $d$ maintains its committed service events
$\mathcal{E}_d$. Writing $t$ for the current decision time, the derived
features are the availability delay
$\max(0,\mathrm{avail}_d-t)$, the number of queued future services
$q_d=|\{e\in\mathcal{E}_d : s_e>t\}|$, the waiting-line saturation
$\min(q_d/W_d,1)$ with $W_d$ the number of physical waiting cells, the
residual time of the service in progress, and the committed workload
$\sum_{e\,\text{unfinished}}\big(c_e-\max(s_e,t)\big)$, together with
the dock class (inbound/outbound) and its coordinates.

\emph{AMR nodes.} Each AMR encodes its busy flag and time-to-free,
per-type inventory ratios and residual capacity, the number of loads
currently onboard, spatial context (coordinates, distance to the
nearest dock, a near-dock indicator), and a local fleet density---the
fraction of other AMRs within a two-cell radius---which is the model's
proxy for imminent interference.

\emph{Job nodes.} Each job encodes its service duration, waiting time
since release, lifecycle status (unpicked/onboard/delivered), carrier
identity, material type, dock indices, and a completion-time estimate
computed with the calibrated travel model
(Section~\ref{sec:calibration}).

\emph{Per-action features.} For every candidate triple $(\omega,k,j)$,
three decision-grade quantities are computed directly: the calibrated
travel time from AMR $k$'s position to the service dock, the calibrated
downstream distance ($o_j$ to $d_j$ for pickups; zero for deliveries),
and the projected queue-aware dock wait $\hat{w}(\cdot)$ of
\eqref{eq:surrogate}. Providing these estimates explicitly relieves the
network from recovering travel arithmetic from coordinates through
several nonlinear layers.

\begin{table}[t]
\caption{Node and Action Features of \textsc{DockGNN}}
\label{tab:features}
\centering
\footnotesize
\begin{tabular}{@{}llc@{}}
\toprule
& Feature & Scale \\
\midrule
\multirow{9}{*}{\rotatebox{90}{AMR $\bm{x}^{A}_k\!\in\!\mathbb{R}^{12}$}}
& busy indicator & -- \\
& time until available & $1/100$ \\
& inventory ratio, per type ($\times3$) & $1/Q$ \\
& residual capacity ratio & $1/(|\mathcal{T}|Q)$ \\
& loads onboard & $1/10$ \\
& near-dock indicator & -- \\
& distance to nearest dock & $1/20$ \\
& local fleet density ($\le\!2$ cells) & $1/(m{-}1)$ \\
& position $(x,y)$ ($\times2$) & $1/10$ \\
\midrule
\multirow{9}{*}{\rotatebox{90}{Job $\bm{x}^{J}_j\!\in\!\mathbb{R}^{11}$}}
& existence flag & -- \\
& service duration $p_j$ & $1/25$ \\
& waiting time since release & $1/100$ \\
& carrier index (0 if unpicked) & $1/m$ \\
& lifecycle status $\{0,\tfrac12,1\}$ & -- \\
& material type one-hot ($\times3$) & -- \\
& inbound dock index & $1/|\mathcal{D}^{\mathrm{in}}|$ \\
& outbound dock index & $1/|\mathcal{D}^{\mathrm{out}}|$ \\
& completion-time estimate & $1/500$ \\
\midrule
\multirow{7}{*}{\rotatebox{90}{Dock $\bm{x}^{D}_d\!\in\!\mathbb{R}^{9}$}}
& inbound / outbound class ($\times2$) & -- \\
& availability delay & $1/100$ \\
& queued services $q_d$ & $1/m$ \\
& waiting-line saturation & -- \\
& residual service in progress & $1/100$ \\
& committed workload & $1/500$ \\
& position $(x,y)$ ($\times2$) & $1/10$ \\
\midrule
\multirow{3}{*}{\rotatebox{90}{Act.\ $\mathbb{R}^{3}$}}
& calibrated travel to service dock & $1/20$ \\
& calibrated downstream distance & $1/20$ \\
& projected queue-aware dock wait & $1/100$ \\
\bottomrule
\end{tabular}
\end{table}

\subsection{Entity Encoders and AMR--Dock Interaction}
\label{sec:interaction}

AMR and dock features pass through two-layer perceptrons and job
features through a linear map,
\begin{equation}
\bm{h}^{A}_k=\mathrm{MLP}_{A}(\bm{x}^{A}_k),\quad
\bm{h}^{D}_d=\mathrm{MLP}_{D}(\bm{x}^{D}_d),\quad
\bar{\bm{h}}^{J}_j=\bm{W}_{\!J}\,\bm{x}^{J}_j,
\end{equation}
all in $\mathbb{R}^{h}$. Fleet--dock coupling is then established by a
single pre-normalization residual attention round
\cite{vaswani2017attention,xiong2020preln} consisting of three
sublayers: docks attend to the fleet, the fleet attends to itself, and
the fleet attends to the updated docks,
\begin{align}
\bm{H}^{D} &\leftarrow \bm{H}^{D} +
\mathrm{MHA}\big(\mathrm{LN}(\bm{H}^{D}),\,\mathrm{LN}(\bm{H}^{A}),\,\mathrm{LN}(\bm{H}^{A})\big),
\label{eq:dock-from-amr}\\
\bm{H}^{A} &\leftarrow \bm{H}^{A} +
\mathrm{MHA}\big(\mathrm{LN}(\bm{H}^{A}),\,\mathrm{LN}(\bm{H}^{A}),\,\mathrm{LN}(\bm{H}^{A})\big),
\label{eq:amr-self}\\
\bm{H}^{A} &\leftarrow \bm{H}^{A} +
\mathrm{MHA}\big(\mathrm{LN}(\bm{H}^{A}),\,\mathrm{LN}(\bm{H}^{D}),\,\mathrm{LN}(\bm{H}^{D})\big),
\label{eq:amr-from-dock}
\end{align}
with four heads per attention. Through
\eqref{eq:dock-from-amr}--\eqref{eq:amr-from-dock}, a dock's embedding
reflects which vehicles are loaded and heading toward it \emph{before}
jobs consume that embedding, and each AMR perceives both the remaining
fleet and network-wide dock congestion.

\emph{Stability.} The output projections of all three attention
sublayers are zero-initialized, making the entire block an exact
identity at the start of training in the spirit of residual-branch
zero initialization \cite{bachlechner2021rezero,goyal2017accurate}.
This is not cosmetic: the embedding stream entering the block has small
magnitude relative to the unit-scale outputs that randomly initialized
attention produces from layer-normalized inputs, and without
zero initialization the three stacked residual additions amplify the
fleet embeddings by more than an order of magnitude at initialization,
which we found to prevent policy-gradient training from converging.

\subsection{Job Contention--Precedence Graph}
\label{sec:graph}

Jobs interact through two mechanisms with different lifetimes:
\emph{future contention} for shared docks, and \emph{committed
precedence} induced by scheduling decisions already made. Both are
encoded in a single dynamic graph $\mathcal{G}_t=(\mathcal{J},E_t)$
with symmetric edges over the job nodes:
(i)~a \emph{contention edge} joins every pair of unfinished jobs that
share an inbound or an outbound dock, so resource competition is
visible from the first decision onward; and
(ii)~\emph{order edges} join consecutively scheduled jobs on the same
AMR and on the same outbound dock, incrementally densifying the graph
as $\sigma_t$ grows.

Before message passing, each job fuses the embeddings of its two docks,
\begin{equation}
\bm{g}^{(0)}_j=\mathrm{MLP}_{F}\!\big(\big[\,\bar{\bm{h}}^{J}_j \,\|\,
\bm{h}^{D}_{o_j} \,\|\, \bm{h}^{D}_{d_j}\big]\big),
\label{eq:fusion}
\end{equation}
so every job representation is conditioned on the live congestion state
of exactly the two resources it will occupy. Message passing then
applies $L=3$ residual GIN layers with \emph{degree-normalized}
aggregation and layer normalization,
\begin{equation}
\bm{g}^{(l+1)}_j=\bm{g}^{(l)}_j+
\mathrm{LN}\!\left(\mathrm{MLP}^{(l)}\!\Big(
(1{+}\epsilon^{(l)})\,\bm{g}^{(l)}_j+
\tfrac{1}{\deg_j}\!\!\sum_{i\in\mathcal{N}(j)}\!\bm{g}^{(l)}_i
\Big)\right)\!.
\label{eq:gin}
\end{equation}
Mean aggregation deliberately trades the injectivity of sum-based GIN
\cite{xu2019gin} for scale invariance: contention edges form dock
cliques whose size grows linearly with $n/|\mathcal{D}|$, and
sum aggregation would make message magnitudes---and hence the
activation statistics of every downstream layer---depend on the
instance size, undermining transfer across job counts. With mean
aggregation the same trained network processes $n{=}10$ and $n{=}100$
instances at identical activation scales.

\subsection{Action Scoring and Value Estimation}
\label{sec:actor}

Let $\bm{e}_{\omega}\in\mathbb{R}^{h}$ be a learned embedding of the
operation type and let $d(\omega,j)=o_j$ if $\omega=\textsc{pick}$ and
$d_j$ otherwise denote the dock at which the candidate action would be
served. The actor scores every triple by
\begin{equation}
z_{\omega,k,j}=\mathrm{MLP}_{\pi}\!\big(\big[\,\bm{e}_{\omega}\,\|\,
\bm{h}^{A}_{k}\,\|\,\bm{g}^{(L)}_{j}\,\|\,
\bm{h}^{D}_{d(\omega,j)}\,\|\,\bm{x}^{\mathrm{act}}_{\omega,k,j}\big]\big),
\label{eq:logit}
\end{equation}
and the policy is the masked softmax
\begin{equation}
\pi_{\theta}(a\,|\,s_t)=
\frac{\exp(z_{a})\,\mathbf{1}[a\in\Lambda(s_t)]}
{\sum_{a'\in\Lambda(s_t)}\exp(z_{a'})}.
\label{eq:policy}
\end{equation}
The critic estimates the state value from pooled entity summaries,
\begin{equation}
V_{\phi}(s_t)=\mathrm{MLP}_{V}\!\big(\big[\,
\overline{\bm{g}}^{\mathrm{(unf)}} \,\|\, \overline{\bm{h}}^{A} \,\|\,
\overline{\bm{h}}^{D}\big]\big),
\label{eq:critic}
\end{equation}
where $\overline{\bm{g}}^{\mathrm{(unf)}}$ averages the embeddings of
unfinished jobs only. A full episode requires exactly $2n$ forward
passes; each pass scores $2mn$ actions in a single batched evaluation
of \eqref{eq:logit}, and all components are permutation-invariant in
the job and AMR sets, so the architecture applies unchanged to any
instance size.

% =========================================================================

\section{Surrogate Model Calibration}
\label{sec:calibration}

\subsection{The Two-Simulator Gap}

The policy constructs schedules on the surrogate $\hat{\Psi}$ of
Section~\ref{sec:mdp}, whereas solution quality is defined by the
executor $\Phi$ of \eqref{eq:objective}. An uncalibrated surrogate that
advances clocks by Manhattan distances and single-scalar dock
availabilities is systematically optimistic: it ignores collision
detours, dock-clearing maneuvers, waiting-line approach paths, and the
upstream holds triggered when waiting lines saturate. The state
features of Section~\ref{sec:features}---and therefore every decision
the policy makes---would then be conditioned on clocks that drift
further from reality as congestion grows, precisely in the regime where
scheduling decisions matter most. Rather than paying the computational
cost of embedding the full executor in the training loop, we close most
of this gap with two lightweight corrections fitted offline.

\subsection{Calibrated Travel and Dock-Wait Estimators}

\emph{Travel.} Realized leg times are modeled as an affine function of
the Manhattan distance, clamped from below by the metric itself:
\begin{equation}
\hat{t}(u,v)=\max\!\big(\Delta(u,v),\ \alpha\,\Delta(u,v)+\beta\big).
\label{eq:travelcal}
\end{equation}

\emph{Dock waiting.} The wait experienced by an AMR that becomes ready
at a dock $d$ at time $\rho$ is modeled as the committed availability
delay plus a penalty indexed by the queue depth at readiness,
\begin{equation}
\hat{w}_d(\rho)=\max\!\big(0,\ \mathrm{avail}_d-\rho\big)
+\varphi\!\big(\min(q_d(\rho),\,3)\big),
\label{eq:waitcal}
\end{equation}
where $q_d(\rho)=|\{e\in\mathcal{E}_d : c_e>\rho\}|$ counts committed
services unfinished at $\rho$, and
$\varphi:\{0,1,2,3^{+}\}\to\mathbb{R}_{\ge0}$ is a monotone lookup
table. The bounded table shape is deliberate: waiting-line detours and
upstream holds produce a cost cliff once the line saturates, a
nonlinearity that a linear congestion model cannot express.

\subsection{Fitting Procedure}

Calibration data are obtained by replaying complete schedules through
both simulators and aligning them per operation. Schedules are drawn
from heterogeneous sources---five dispatching-rule pairs and the
current learned policies---across instance sizes
$n\in\{10,20,40,60\}$, so that the fitted correction covers the state
distribution that policies actually visit. For every operation, the
raw surrogate supplies the Manhattan leg distance, the base wait, and
the queue depth at readiness; the executor's logged timeline is
partitioned, per AMR, into legs terminating at each service start,
within which movement entries sum to the realized travel and the
remainder is attributed to waiting. Executions containing infeasible
operations are discarded. The travel parameters $(\alpha,\beta)$ are
fitted by least squares and the penalties $\varphi(\cdot)$ as clipped
bin means with monotonicity enforced across depths.

In our environment, fitting on $10{,}360$ operation pairs from $159$
schedules yielded $\alpha=0.963$, $\beta=1.12$ ($R^{2}=0.91$) and
$\varphi=(0.04,\,0.04,\,1.66,\,3.57)$; the step between depths one and
two locates the waiting-line saturation cliff. The correction removes
the mean per-leg travel bias entirely ($+0.744\to0.000$ time units)
and halves the end-to-end makespan drift between surrogate and
executor (mean absolute error $24.3\to12.1$ on held-out schedules).
The calibration is a frozen artifact: it is fitted once before
training, adds no learned components to the training loop, and is
shared by \emph{all} policies and by the dispatching-rule baseline of
Section~\ref{sec:training}, so every method plans in the same
corrected world. After training converges, the fit can be re-estimated
from the new policies' schedules to confirm stability.

% =========================================================================

\section{Policy Optimization}
\label{sec:training}

\subsection{Objective and Gradient Estimator}

The policy $\pi_\theta$ is trained to minimize the expected executed
makespan $\mathbb{E}_{\pi_\theta}[C_{\max}(\Phi(\bm{a},\sigma))]$ with
the REINFORCE estimator \cite{williams1992reinforce}. During training,
episodes are generated by sampling from the masked policy
\eqref{eq:policy} on the calibrated surrogate; at validation and
deployment, actions are selected greedily. The gradient quality of
REINFORCE hinges on its baseline; we use a counterfactual baseline
built from a strong dispatching rule and evaluated by the executor
itself.

\subsection{Counterfactual Dispatch-Rule Baseline}
\label{sec:baseline}

Let $R$ denote a deterministic dispatching heuristic (job rule and AMR
rule; we use earliest-completion for both) and let
$R(\sigma_{1:t})$ denote the complete solution obtained by executing
the partial sequence $\sigma_{1:t}$ and completing all remaining
operations with $R$. Define the shorthand
$C[\cdot]=C_{\max}(\Phi(\cdot))$. The advantage of the action taken at
step $t$ is the improvement of committing that action and then
following the rule, relative to following the rule immediately:
\begin{equation}
A_t \;=\; C\big[R(\sigma_{1:t-1})\big]\;-\;
C\big[R(\sigma_{1:t-1}\oplus a_t)\big].
\label{eq:advantage}
\end{equation}
Both completions are evaluated by the collision-aware executor, so the
learning signal is grounded in true solution quality rather than in
surrogate clocks. Unlike episode-level rollout baselines
\cite{rennie2017self,kool2019attention}, \eqref{eq:advantage} assigns
credit per decision: an early pickup that congests a dock is penalized
at the step where it is committed, not diluted across the episode.
The price is $\mathcal{O}(n)$ executor evaluations per episode, which
dominates training cost; we consider the sharper credit assignment a
favorable trade at the instance sizes studied here, and
Section~\ref{sec:training-details} lists a cheaper actor--critic
alternative.

\subsection{Loss Function}

Advantages are standardized to zero mean and unit variance across all
steps of a batch of $B$ episodes. The per-batch loss combines the
policy gradient with a load-balancing shaping term and an entropy
bonus:
\begin{equation}
\mathcal{L}(\theta)=
-\frac{1}{B}\sum_{b=1}^{B}\sum_{t=1}^{2n}
\Big[\log\pi_\theta(a_t^b|s_t^b)\,
\big(\hat{A}_t^{\,b}+\lambda_{\mathrm{lb}}A_t^{\mathrm{lb},b}\big)
+\lambda_{H}\,\mathcal{H}_t^{b}\Big],
\label{eq:loss}
\end{equation}
where $\hat{A}_t$ is the standardized advantage,
$\mathcal{H}_t$ is the entropy of the masked policy over feasible
actions, and the shaping term
$A^{\mathrm{lb}}_t=\min_{k'}c_{k'}-c_{k}$ (for pickups; zero for
deliveries) compares the assignment count $c_k$ of the chosen AMR with
the least-loaded AMR, discouraging degenerate fleet utilization early
in training \cite{ng1999shaping}. We use
$\lambda_{\mathrm{lb}}=0.1$ and $\lambda_{H}=0.01$.

\begin{algorithm}[t]
\caption{Training with counterfactual dispatch baseline}
\label{alg:training}
\begin{algorithmic}[1]
\Require policy $\pi_\theta$, rule $R$, executor $\Phi$, surrogate $\hat{\Psi}$
\State fit calibration $(\alpha,\beta,\varphi)$ offline; freeze (Sec.~\ref{sec:calibration})
\For{$\mathrm{epoch}=1,\dots,E$}
  \For{$b=1,\dots,B$}
    \State sample instance; roll out $\sigma^b\!\sim\!\pi_\theta$ on $\hat{\Psi}$
    \For{$t=1,\dots,2n$}
      \State $A^b_t \gets C[R(\sigma^b_{1:t-1})]-C[R(\sigma^b_{1:t-1}\oplus a^b_t)]$
    \EndFor
  \EndFor
  \State standardize $\{A^b_t\}$ over the batch
  \State update $\theta$ by Adam on $\mathcal{L}(\theta)$ in \eqref{eq:loss}; clip $\lVert\nabla\rVert\le1$
  \If{$\mathrm{epoch}\bmod T_{\mathrm{val}}=0$}
    \State greedy-decode held-out instances; score by
    $C_{\max}+\kappa\cdot(\text{infeasible ops})$; keep best
  \EndIf
\EndFor
\end{algorithmic}
\end{algorithm}

\subsection{Training Protocol}
\label{sec:training-details}

Algorithm~\ref{alg:training} summarizes the procedure. Only the actor
parameters (encoders, interaction block, GIN layers, operation
embedding, and actor head) receive gradients under REINFORCE; the
critic \eqref{eq:critic} is reserved for an optional
proximal-policy-optimization variant \cite{schulman2017ppo} that
replaces \eqref{eq:advantage} with critic-based advantages when the
executor evaluations of the stepwise baseline are prohibitive. We
train with Adam \cite{kingma2015adam} (learning rate $3\times10^{-4}$,
batch $B=16$, gradient-norm clip $1.0$) for $2{,}000$ epochs. Every
$T_{\mathrm{val}}=50$ epochs the policy is evaluated by greedy decoding on a fixed
held-out validation set, disjoint from the training pool, and scored
by the executed makespan plus a penalty of $\kappa=10^{3}$ per
infeasible operation; the checkpoint with the best validation score is
retained. Two safeguards proved important in practice: the
identity-at-initialization residual branches of
Section~\ref{sec:interaction}, without which policy-gradient training
failed to converge, and computing the entropy in \eqref{eq:loss} over
the feasible-action support only, which keeps the bonus comparable as
the mask tightens toward the end of an episode.

% =========================================================================
% Suggested BibTeX entries (replace keys above with your bibliography):
%
% @inproceedings{vaswani2017attention, Vaswani et al., "Attention Is All
%   You Need", NeurIPS 2017}
% @inproceedings{xu2019gin, Xu, Hu, Leskovec, Jegelka, "How Powerful Are
%   Graph Neural Networks?", ICLR 2019}
% @inproceedings{xiong2020preln, Xiong et al., "On Layer Normalization in
%   the Transformer Architecture", ICML 2020}
% @inproceedings{bachlechner2021rezero, Bachlechner et al., "ReZero Is
%   All You Need: Fast Convergence at Large Depth", UAI 2021}
% @article{goyal2017accurate, Goyal et al., "Accurate, Large Minibatch
%   SGD: Training ImageNet in 1 Hour", arXiv:1706.02677}
% @inproceedings{zhang2020l2d, Zhang et al., "Learning to Dispatch for
%   Job Shop Scheduling via Deep Reinforcement Learning", NeurIPS 2020}
% @article{song2023fjsp, Song et al., "Flexible Job-Shop Scheduling via
%   Graph Neural Network and Deep Reinforcement Learning", IEEE TII 2023}
% @inproceedings{stern2019mapf, Stern et al., "Multi-Agent Pathfinding:
%   Definitions, Variants, and Benchmarks", SoCS 2019}
% @article{williams1992reinforce, Williams, "Simple Statistical
%   Gradient-Following Algorithms for Connectionist Reinforcement
%   Learning", Machine Learning 8, 1992}
% @inproceedings{rennie2017self, Rennie et al., "Self-Critical Sequence
%   Training for Image Captioning", CVPR 2017}
% @inproceedings{kool2019attention, Kool, van Hoof, Welling, "Attention,
%   Learn to Solve Routing Problems!", ICLR 2019}
% @inproceedings{ng1999shaping, Ng, Harada, Russell, "Policy Invariance
%   Under Reward Transformations: Theory and Application to Reward
%   Shaping", ICML 1999}
% @article{schulman2017ppo, Schulman et al., "Proximal Policy
%   Optimization Algorithms", arXiv:1707.06347}
% @inproceedings{kingma2015adam, Kingma, Ba, "Adam: A Method for
%   Stochastic Optimization", ICLR 2015}
% =========================================================================
 inside an IEEEtran document.
% Requires in the preamble:
%   \usepackage{amsmath,amssymb,bm}
%   \usepackage{booktabs,array,multirow,graphicx}
%   \usepackage{algorithm}
%   \usepackage{algpseudocode}
% Citation keys are placeholders -- suggested BibTeX entries are listed in
% comments at the end of this file. Replace \cite{...} keys with your own.
% "DockGNN" is a placeholder model name; rename globally if you prefer.
% Contains four sections: Problem Formulation, Architecture, Surrogate
% Calibration (sec:calibration), and Policy Optimization (sec:training).
% =========================================================================

\section{Problem Formulation}
\label{sec:problem}

\begin{table*}[t]
\caption{Summary of Notation}
\label{tab:notation}
\centering
\footnotesize
\begin{tabular}{@{}l p{0.295\textwidth} @{\qquad} l p{0.295\textwidth}@{}}
\toprule
Symbol & Meaning & Symbol & Meaning \\
\midrule
\multicolumn{2}{@{}l}{\itshape Problem setting} &
\multicolumn{2}{@{}l}{\itshape Architecture} \\
$\mathcal{W}$ & grid workspace (four-connected unit cells) &
  $\bm{x}^{A}_k,\bm{x}^{J}_j,\bm{x}^{D}_d$ & AMR / job / dock input features ($\mathbb{R}^{12},\mathbb{R}^{11},\mathbb{R}^{9}$) \\
$\Delta(u,v)$ & Manhattan distance between positions $u,v$ &
  $\bm{x}^{\mathrm{act}}_{\omega,k,j}$ & per-action decision features ($\mathbb{R}^{3}$) \\
$\mathcal{J},\ n$ & job set; number of jobs &
  $h$ & hidden width of all embeddings ($h{=}128$) \\
$\mathcal{M},\ m$ & AMR fleet; fleet size &
  $\bm{h}^{A}_k,\ \bm{h}^{D}_d$ & AMR and dock embeddings \\
$\mathcal{T}$ & set of material types &
  $\bar{\bm{h}}^{J}_j$ & job embedding before dock fusion \\
$\mathcal{D}^{\mathrm{in}},\mathcal{D}^{\mathrm{out}},\mathcal{D}$ & inbound docks; outbound docks; their union &
  $\bm{H}^{A},\bm{H}^{D}$ & stacked AMR / dock embedding matrices \\
$\tau_j$ & material type of job $j$ &
  $\mathrm{MHA},\mathrm{LN}$ & multi-head attention; layer normalization \\
$o_j,\ d_j$ & inbound and outbound dock of job $j$ &
  $\bm{g}^{(l)}_j$ & job embedding after fusion and $l$ GIN layers \\
$r_j$ & release time of job $j$ &
  $\mathcal{G}_t=(\mathcal{J},E_t)$ & contention--precedence graph at step $t$ \\
$p_j=p(\tau_j)$ & dock service duration of job $j$ &
  $\mathcal{N}(j),\ \deg_j$ & graph neighbors and degree of job $j$ \\
$\pi_j,\ \delta_j$ & pickup and delivery operations of job $j$ &
  $\epsilon^{(l)},\ L$ & learnable GIN self-weight; number of GIN layers \\
$Q$ & per-type AMR carrying capacity &
  $\bm{e}_{\omega}$ & learned operation-type embedding \\
$\bm{a},\ a_j$ & AMR assignment vector; AMR assigned to job $j$ &
  $d(\omega,j)$ & dock served by the action ($o_j$ if pickup, else $d_j$) \\
$\sigma,\ \sigma_t$ & operation sequence; prefix committed up to step $t$ &
  $z_{\omega,k,j}$ & logit of action $(\omega,k,j)$ \\
$s^{\pi}_j,c^{\pi}_j$ & start / completion of pickup service of job $j$ &
  $V_{\phi}$ & critic (state-value) network, parameters $\phi$ \\
$s^{\delta}_j,c^{\delta}_j$ & start / completion of delivery service of job $j$ &
  & \\
$t_k(u,v)$ & realized travel time of AMR $k$ from $u$ to $v$ &
  \multicolumn{2}{@{}l}{\itshape Training} \\
$\Phi$ & collision-aware executor &
  $R,\ R(\sigma_{1:t})$ & dispatching rule; completion of prefix $\sigma_{1:t}$ by $R$ \\
$C_{\max},\ F_k$ & makespan; finish time of AMR $k$ incl.\ depot return &
  $C[\cdot]$ & shorthand for $C_{\max}(\Phi(\cdot))$ \\
$W_d$ & number of physical waiting cells at dock $d$ &
  $A_t,\ \hat{A}_t$ & stepwise counterfactual advantage; standardized \\
\multicolumn{2}{@{}l}{\itshape Sequential decision process} &
  $A^{\mathrm{lb}}_t$ & load-balance shaping advantage \\
$s_t,\ a_t$ & state and action at decision epoch $t$ &
  $\lambda_{\mathrm{lb}},\ \lambda_{H}$ & shaping and entropy coefficients \\
$\omega$ & operation type, $\omega\in\{\textsc{pick},\textsc{deliver}\}$ &
  $\mathcal{H}_t$ & entropy of the masked policy at step $t$ \\
$\Lambda(s_t)$ & feasible-action set (mask) in state $s_t$ &
  $\mathcal{L}(\theta)$ & training loss \\
$\hat{\Psi}$ & calibrated surrogate transition model &
  $B,\ E$ & batch size (episodes); number of training epochs \\
$\mathcal{E}_d$ & committed service events $(s_e,c_e,j_e,\omega_e)$ at dock $d$ &
  $T_{\mathrm{val}},\ \kappa$ & validation interval; infeasibility penalty \\
$\hat{b}_k,\ x_k$ & next-available time and position of AMR $k$ &
  $\mathbf{1}[\cdot]$ & indicator function \\
$\rho$ & ready time of an AMR at a dock &
  & \\
$\hat{t}(u,v)$ & calibrated travel-time estimate \eqref{eq:travelcal} &
  & \\
$\hat{w}_d(\rho)$ & queue-aware dock-wait estimate \eqref{eq:waitcal} &
  & \\
$q_d(\rho)$ & services still unfinished at time $\rho$ at dock $d$ &
  & \\
$\alpha,\ \beta$ & travel-calibration scale and offset &
  & \\
$\varphi(\cdot)$ & dock-wait penalty table over queue depths &
  & \\
$\pi_{\theta}$ & scheduling policy with parameters $\theta$ &
  & \\
\bottomrule
\end{tabular}
\end{table*}

\subsection{System Description}

We consider the internal transfer problem of a cross-docking terminal
served by a fleet of autonomous mobile robots (AMRs). Table~\ref{tab:notation}
summarizes the notation used throughout the paper. The terminal floor
is modeled as a four-connected unit grid $\mathcal{W}$ in which a set of
\emph{inbound docks} $\mathcal{D}^{\mathrm{in}}$ is located along one
boundary and a set of \emph{outbound docks} (processing stations)
$\mathcal{D}^{\mathrm{out}}$ along the opposite boundary; the AMR depots
lie between them. Consistent with cross-docking practice, the floor is
free of static obstacles, so congestion arises from vehicle interference
and dock contention rather than from the layout itself. AMRs travel at
unit speed between adjacent cells, and the distance between two
positions $u,v$ is the Manhattan metric $\Delta(u,v)$.

A set of transfer jobs $\mathcal{J}=\{1,\dots,n\}$ must be moved across
the terminal by a fleet $\mathcal{M}=\{1,\dots,m\}$ of homogeneous AMRs.
Each job $j\in\mathcal{J}$ represents one unit load of a material type
$\tau_j\in\mathcal{T}$ that becomes available at its inbound dock
$o_j\in\mathcal{D}^{\mathrm{in}}$ at release time $r_j\ge 0$ and must be
delivered to a designated outbound dock
$d_j\in\mathcal{D}^{\mathrm{out}}$. Handling a unit load of type $\tau$
occupies a dock for a type-dependent service duration $p(\tau)$; we
write $p_j=p(\tau_j)$ and assume the loading service at the inbound dock
and the unloading/processing service at the outbound dock take the same
time, reflecting symmetric handling of a unit load. Each job therefore
consists of an ordered pair of operations: a \emph{pickup} operation
$\pi_j$ executed at $o_j$ and a \emph{delivery} operation $\delta_j$
executed at $d_j$, both performed by the same AMR (no transshipment).
An AMR may consolidate loads subject to a per-type carrying capacity
$Q$: at no time may it carry more than $Q$ units of any single material
type.

Docks are exclusive single-server resources: at most one AMR may be in
service at a dock at any time, and services are non-preemptive. An AMR
arriving at an occupied dock queues in a bounded set of waiting cells
adjacent to that dock; when the waiting line is saturated, the AMR must
hold at an upstream position and re-approach later. Finally, AMR motion
must be collision-free: at any unit time step, each grid cell may be
occupied by at most one AMR.

\subsection{Formal Statement}

A solution is a pair $(\bm{a},\sigma)$, where
$\bm{a}=(a_1,\dots,a_n)\in\mathcal{M}^{n}$ assigns each job to an AMR
and $\sigma$ is a total order over the $2n$ operations
$\{\pi_j,\delta_j\}_{j\in\mathcal{J}}$ that determines the sequence in
which operations are committed. Let $s^{\pi}_j,c^{\pi}_j$
(resp.\ $s^{\delta}_j,c^{\delta}_j$) denote the service start and
completion times of $\pi_j$ (resp.\ $\delta_j$) induced by executing
$(\bm{a},\sigma)$. A feasible execution satisfies:
\begin{align}
c^{\pi}_j &= s^{\pi}_j + p_j, \qquad
c^{\delta}_j = s^{\delta}_j + p_j,
&& \forall j\in\mathcal{J}, \label{eq:nonpreempt}\\
s^{\pi}_j &\ge r_j,
&& \forall j\in\mathcal{J}, \label{eq:release}\\
s^{\delta}_j &\ge c^{\pi}_j,
&& \forall j\in\mathcal{J}, \label{eq:precedence}
\end{align}
\begin{align}
\big|\{\,j : a_j{=}k,\ \tau_j{=}\tau,\ c^{\pi}_j \le u < c^{\delta}_j \}\big| &\le Q,
&& \forall k,\tau,u, \label{eq:capacity}\\
\big[s_e, c_e\big) \cap \big[s_{e'}, c_{e'}\big) &= \emptyset,
&& \forall e\ne e' \ \text{at the same dock}, \label{eq:exclusive}
\end{align}
where \eqref{eq:nonpreempt} enforces non-preemptive services,
\eqref{eq:release} release dates, \eqref{eq:precedence} pickup-before-%
delivery precedence, \eqref{eq:capacity} the per-type carrying capacity,
and \eqref{eq:exclusive} dock exclusivity over all service events $e$.
In addition, for any two operations $u\prec v$ executed consecutively by
the same AMR $k$,
\begin{equation}
s_v \;\ge\; c_u + t_k\big(\mathrm{loc}(u),\mathrm{loc}(v)\big),
\label{eq:travel}
\end{equation}
where $t_k(\cdot,\cdot)\ge\Delta(\cdot,\cdot)$ is the realized travel
time of AMR $k$, which may exceed the Manhattan distance due to
collision avoidance, queuing detours, and dock-clearing maneuvers.

The realized travel and waiting times are not fixed constants but are
determined by a high-fidelity \emph{executor} $\Phi$, which routes every
AMR leg with space--time $\mathrm{A}^{*}$ search over a shared
reservation table (at most one AMR per cell per time step), manages the
bounded dock waiting lines, and inserts upstream holding maneuvers when
lines are saturated. The problem is therefore a simulation-based
optimization: find
\begin{equation}
(\bm{a}^{*},\sigma^{*}) \;=\;
\arg\min_{(\bm{a},\sigma)} \; C_{\max}\big(\Phi(\bm{a},\sigma)\big),
\label{eq:objective}
\end{equation}
where the makespan $C_{\max}=\max_{k\in\mathcal{M}} F_k$ and $F_k$ is
the time at which AMR $k$ completes its final activity, including its
return trip to the depot. The problem generalizes both the capacitated
pickup-and-delivery problem and parallel-machine scheduling with
sequence-dependent setups, and is therefore NP-hard; the coupling to
collision-free multi-agent routing further compounds its difficulty
\cite{stern2019mapf}.

\subsection{Sequential Decision Formulation}
\label{sec:mdp}

We cast the construction of $(\bm{a},\sigma)$ as an episodic sequential
decision process with exactly $2n$ decision epochs, one per committed
operation. At epoch $t$, the state
\begin{equation}
s_t=\big(\{\bm{x}^{A}_{k}\},\, \{\bm{x}^{J}_{j}\},\, \{\bm{x}^{D}_{d}\},\,
\{\mathcal{E}_d\},\, \sigma_{t}\big)
\label{eq:state}
\end{equation}
collects the AMR states (position, next-available time, per-type
inventory), the job states (unpicked / onboard / delivered, carrier
identity), the dock states, the set of committed dock service events
$\mathcal{E}_d=\{(s_e,c_e,j_e,\omega_e)\}$ at each dock $d$, and the
partial operation sequence $\sigma_t$.

An action is a triple
\begin{equation}
a_t=(\omega,k,j)\in\{\textsc{pick},\textsc{deliver}\}\times\mathcal{M}\times\mathcal{J},
\label{eq:action}
\end{equation}
giving a flat action space of size $2mn$ restricted by a feasibility
mask $\Lambda(s_t)$: $(\textsc{pick},k,j)$ is legal iff $j$ is unpicked
and AMR $k$ carries fewer than $Q$ units of type $\tau_j$;
$(\textsc{deliver},k,j)$ is legal iff $j$ is onboard AMR $k$. The mask
guarantees by construction that every complete episode yields a feasible
solution, so no repair mechanism is required.

Transitions follow a deterministic surrogate model $\hat{\Psi}$ that
advances the acting AMR's clock using a \emph{calibrated} travel-time
estimate $\hat{t}(\cdot,\cdot)$ and a queue-aware dock-wait estimate
$\hat{w}_d(\cdot)$ fitted offline against the executor $\Phi$
(Section~\ref{sec:calibration}); committing an action appends the
corresponding service window to $\mathcal{E}_d$. For a pickup, the
service start of job $j$ by AMR $k$ located at $x_k$ with next-available
time $\hat{b}_k$ is
\begin{equation}
s^{\pi}_j=\underbrace{\max\!\big(\hat{b}_k+\hat{t}(x_k,o_j),\, r_j\big)}_{\text{ready time }\rho}
\;+\;\hat{w}_{o_j}(\rho),
\label{eq:surrogate}
\end{equation}
and analogously for deliveries at $d_j$. After $2n$ steps the completed
solution is evaluated by the executor, and the learning objective is to
minimize the expected executed makespan
$\mathbb{E}_{\pi_\theta}\!\left[C_{\max}(\Phi(\bm{a},\sigma))\right]$;
the policy-gradient training procedure and its dispatch-rule baseline
are described in Section~\ref{sec:training}.

% =========================================================================

\section{Dock-Aware Graph Scheduler}
\label{sec:architecture}

\subsection{Overview and Design Rationale}

Existing learning-to-dispatch architectures for shop scheduling encode
jobs (operations) as graph nodes and machines as flat feature vectors
\cite{zhang2020l2d,song2023fjsp}. Applied to cross-docking, this design
has three blind spots. First, the binding resources are the docks: with
$m$ AMRs contending for $|\mathcal{D}^{\mathrm{in}}|+|\mathcal{D}^{\mathrm{out}}|$
single-server docks, queuing at docks---not processing---dominates the
makespan, yet job-only graphs represent dock congestion at best as
scalar features replicated across jobs. Second, the temporal structure
of dock occupancy (which service windows are already committed, and
when) is invisible. Third, the AMR fleet is encoded independently
per vehicle, so mutual interference cannot be anticipated.

The proposed scheduler, hereafter \textsc{DockGNN}, addresses all three
by treating docks as first-class graph entities. It comprises
(i)~heterogeneous encoders for AMR, job, and dock nodes;
(ii)~an AMR--dock interaction block that exchanges information between
the fleet and the dock set;
(iii)~a job \emph{contention--precedence} graph processed by
degree-normalized graph isomorphism network (GIN) layers
\cite{xu2019gin}; and
(iv)~a factored actor head that scores each masked action from the
operation type, the acting AMR, the target job, the operation-specific
dock, and a small set of per-action decision features. The full model
has $6.2\times10^{5}$ parameters ($5.7\times10^{5}$ actor,
$4.9\times10^{4}$ critic) with hidden width $h=128$ throughout.

% TODO: architecture overview figure
% \begin{figure*}[t]
%   \centering
%   \includegraphics[width=\textwidth]{figures/dockgnn_architecture}
%   \caption{Overview of \textsc{DockGNN}. Dock, AMR, and job features are
%   encoded separately; an AMR--dock attention block exchanges congestion
%   information; jobs fuse their inbound/outbound dock embeddings and pass
%   messages over the contention--precedence graph; the actor scores each
%   feasible (operation, AMR, job) triple.}
%   \label{fig:architecture}
% \end{figure*}

\subsection{Heterogeneous State Representation}
\label{sec:features}

At every decision epoch the state \eqref{eq:state} is rendered into four
feature families, summarized in Table~\ref{tab:features}. All features
are normalized to comparable magnitudes by fixed scales.

\emph{Dock nodes.} Each dock $d$ maintains its committed service events
$\mathcal{E}_d$. Writing $t$ for the current decision time, the derived
features are the availability delay
$\max(0,\mathrm{avail}_d-t)$, the number of queued future services
$q_d=|\{e\in\mathcal{E}_d : s_e>t\}|$, the waiting-line saturation
$\min(q_d/W_d,1)$ with $W_d$ the number of physical waiting cells, the
residual time of the service in progress, and the committed workload
$\sum_{e\,\text{unfinished}}\big(c_e-\max(s_e,t)\big)$, together with
the dock class (inbound/outbound) and its coordinates.

\emph{AMR nodes.} Each AMR encodes its busy flag and time-to-free,
per-type inventory ratios and residual capacity, the number of loads
currently onboard, spatial context (coordinates, distance to the
nearest dock, a near-dock indicator), and a local fleet density---the
fraction of other AMRs within a two-cell radius---which is the model's
proxy for imminent interference.

\emph{Job nodes.} Each job encodes its service duration, waiting time
since release, lifecycle status (unpicked/onboard/delivered), carrier
identity, material type, dock indices, and a completion-time estimate
computed with the calibrated travel model
(Section~\ref{sec:calibration}).

\emph{Per-action features.} For every candidate triple $(\omega,k,j)$,
three decision-grade quantities are computed directly: the calibrated
travel time from AMR $k$'s position to the service dock, the calibrated
downstream distance ($o_j$ to $d_j$ for pickups; zero for deliveries),
and the projected queue-aware dock wait $\hat{w}(\cdot)$ of
\eqref{eq:surrogate}. Providing these estimates explicitly relieves the
network from recovering travel arithmetic from coordinates through
several nonlinear layers.

\begin{table}[t]
\caption{Node and Action Features of \textsc{DockGNN}}
\label{tab:features}
\centering
\footnotesize
\begin{tabular}{@{}llc@{}}
\toprule
& Feature & Scale \\
\midrule
\multirow{9}{*}{\rotatebox{90}{AMR $\bm{x}^{A}_k\!\in\!\mathbb{R}^{12}$}}
& busy indicator & -- \\
& time until available & $1/100$ \\
& inventory ratio, per type ($\times3$) & $1/Q$ \\
& residual capacity ratio & $1/(|\mathcal{T}|Q)$ \\
& loads onboard & $1/10$ \\
& near-dock indicator & -- \\
& distance to nearest dock & $1/20$ \\
& local fleet density ($\le\!2$ cells) & $1/(m{-}1)$ \\
& position $(x,y)$ ($\times2$) & $1/10$ \\
\midrule
\multirow{9}{*}{\rotatebox{90}{Job $\bm{x}^{J}_j\!\in\!\mathbb{R}^{11}$}}
& existence flag & -- \\
& service duration $p_j$ & $1/25$ \\
& waiting time since release & $1/100$ \\
& carrier index (0 if unpicked) & $1/m$ \\
& lifecycle status $\{0,\tfrac12,1\}$ & -- \\
& material type one-hot ($\times3$) & -- \\
& inbound dock index & $1/|\mathcal{D}^{\mathrm{in}}|$ \\
& outbound dock index & $1/|\mathcal{D}^{\mathrm{out}}|$ \\
& completion-time estimate & $1/500$ \\
\midrule
\multirow{7}{*}{\rotatebox{90}{Dock $\bm{x}^{D}_d\!\in\!\mathbb{R}^{9}$}}
& inbound / outbound class ($\times2$) & -- \\
& availability delay & $1/100$ \\
& queued services $q_d$ & $1/m$ \\
& waiting-line saturation & -- \\
& residual service in progress & $1/100$ \\
& committed workload & $1/500$ \\
& position $(x,y)$ ($\times2$) & $1/10$ \\
\midrule
\multirow{3}{*}{\rotatebox{90}{Act.\ $\mathbb{R}^{3}$}}
& calibrated travel to service dock & $1/20$ \\
& calibrated downstream distance & $1/20$ \\
& projected queue-aware dock wait & $1/100$ \\
\bottomrule
\end{tabular}
\end{table}

\subsection{Entity Encoders and AMR--Dock Interaction}
\label{sec:interaction}

AMR and dock features pass through two-layer perceptrons and job
features through a linear map,
\begin{equation}
\bm{h}^{A}_k=\mathrm{MLP}_{A}(\bm{x}^{A}_k),\quad
\bm{h}^{D}_d=\mathrm{MLP}_{D}(\bm{x}^{D}_d),\quad
\bar{\bm{h}}^{J}_j=\bm{W}_{\!J}\,\bm{x}^{J}_j,
\end{equation}
all in $\mathbb{R}^{h}$. Fleet--dock coupling is then established by a
single pre-normalization residual attention round
\cite{vaswani2017attention,xiong2020preln} consisting of three
sublayers: docks attend to the fleet, the fleet attends to itself, and
the fleet attends to the updated docks,
\begin{align}
\bm{H}^{D} &\leftarrow \bm{H}^{D} +
\mathrm{MHA}\big(\mathrm{LN}(\bm{H}^{D}),\,\mathrm{LN}(\bm{H}^{A}),\,\mathrm{LN}(\bm{H}^{A})\big),
\label{eq:dock-from-amr}\\
\bm{H}^{A} &\leftarrow \bm{H}^{A} +
\mathrm{MHA}\big(\mathrm{LN}(\bm{H}^{A}),\,\mathrm{LN}(\bm{H}^{A}),\,\mathrm{LN}(\bm{H}^{A})\big),
\label{eq:amr-self}\\
\bm{H}^{A} &\leftarrow \bm{H}^{A} +
\mathrm{MHA}\big(\mathrm{LN}(\bm{H}^{A}),\,\mathrm{LN}(\bm{H}^{D}),\,\mathrm{LN}(\bm{H}^{D})\big),
\label{eq:amr-from-dock}
\end{align}
with four heads per attention. Through
\eqref{eq:dock-from-amr}--\eqref{eq:amr-from-dock}, a dock's embedding
reflects which vehicles are loaded and heading toward it \emph{before}
jobs consume that embedding, and each AMR perceives both the remaining
fleet and network-wide dock congestion.

\emph{Stability.} The output projections of all three attention
sublayers are zero-initialized, making the entire block an exact
identity at the start of training in the spirit of residual-branch
zero initialization \cite{bachlechner2021rezero,goyal2017accurate}.
This is not cosmetic: the embedding stream entering the block has small
magnitude relative to the unit-scale outputs that randomly initialized
attention produces from layer-normalized inputs, and without
zero initialization the three stacked residual additions amplify the
fleet embeddings by more than an order of magnitude at initialization,
which we found to prevent policy-gradient training from converging.

\subsection{Job Contention--Precedence Graph}
\label{sec:graph}

Jobs interact through two mechanisms with different lifetimes:
\emph{future contention} for shared docks, and \emph{committed
precedence} induced by scheduling decisions already made. Both are
encoded in a single dynamic graph $\mathcal{G}_t=(\mathcal{J},E_t)$
with symmetric edges over the job nodes:
(i)~a \emph{contention edge} joins every pair of unfinished jobs that
share an inbound or an outbound dock, so resource competition is
visible from the first decision onward; and
(ii)~\emph{order edges} join consecutively scheduled jobs on the same
AMR and on the same outbound dock, incrementally densifying the graph
as $\sigma_t$ grows.

Before message passing, each job fuses the embeddings of its two docks,
\begin{equation}
\bm{g}^{(0)}_j=\mathrm{MLP}_{F}\!\big(\big[\,\bar{\bm{h}}^{J}_j \,\|\,
\bm{h}^{D}_{o_j} \,\|\, \bm{h}^{D}_{d_j}\big]\big),
\label{eq:fusion}
\end{equation}
so every job representation is conditioned on the live congestion state
of exactly the two resources it will occupy. Message passing then
applies $L=3$ residual GIN layers with \emph{degree-normalized}
aggregation and layer normalization,
\begin{equation}
\bm{g}^{(l+1)}_j=\bm{g}^{(l)}_j+
\mathrm{LN}\!\left(\mathrm{MLP}^{(l)}\!\Big(
(1{+}\epsilon^{(l)})\,\bm{g}^{(l)}_j+
\tfrac{1}{\deg_j}\!\!\sum_{i\in\mathcal{N}(j)}\!\bm{g}^{(l)}_i
\Big)\right)\!.
\label{eq:gin}
\end{equation}
Mean aggregation deliberately trades the injectivity of sum-based GIN
\cite{xu2019gin} for scale invariance: contention edges form dock
cliques whose size grows linearly with $n/|\mathcal{D}|$, and
sum aggregation would make message magnitudes---and hence the
activation statistics of every downstream layer---depend on the
instance size, undermining transfer across job counts. With mean
aggregation the same trained network processes $n{=}10$ and $n{=}100$
instances at identical activation scales.

\subsection{Action Scoring and Value Estimation}
\label{sec:actor}

Let $\bm{e}_{\omega}\in\mathbb{R}^{h}$ be a learned embedding of the
operation type and let $d(\omega,j)=o_j$ if $\omega=\textsc{pick}$ and
$d_j$ otherwise denote the dock at which the candidate action would be
served. The actor scores every triple by
\begin{equation}
z_{\omega,k,j}=\mathrm{MLP}_{\pi}\!\big(\big[\,\bm{e}_{\omega}\,\|\,
\bm{h}^{A}_{k}\,\|\,\bm{g}^{(L)}_{j}\,\|\,
\bm{h}^{D}_{d(\omega,j)}\,\|\,\bm{x}^{\mathrm{act}}_{\omega,k,j}\big]\big),
\label{eq:logit}
\end{equation}
and the policy is the masked softmax
\begin{equation}
\pi_{\theta}(a\,|\,s_t)=
\frac{\exp(z_{a})\,\mathbf{1}[a\in\Lambda(s_t)]}
{\sum_{a'\in\Lambda(s_t)}\exp(z_{a'})}.
\label{eq:policy}
\end{equation}
The critic estimates the state value from pooled entity summaries,
\begin{equation}
V_{\phi}(s_t)=\mathrm{MLP}_{V}\!\big(\big[\,
\overline{\bm{g}}^{\mathrm{(unf)}} \,\|\, \overline{\bm{h}}^{A} \,\|\,
\overline{\bm{h}}^{D}\big]\big),
\label{eq:critic}
\end{equation}
where $\overline{\bm{g}}^{\mathrm{(unf)}}$ averages the embeddings of
unfinished jobs only. A full episode requires exactly $2n$ forward
passes; each pass scores $2mn$ actions in a single batched evaluation
of \eqref{eq:logit}, and all components are permutation-invariant in
the job and AMR sets, so the architecture applies unchanged to any
instance size.

% =========================================================================

\section{Surrogate Model Calibration}
\label{sec:calibration}

\subsection{The Two-Simulator Gap}

The policy constructs schedules on the surrogate $\hat{\Psi}$ of
Section~\ref{sec:mdp}, whereas solution quality is defined by the
executor $\Phi$ of \eqref{eq:objective}. An uncalibrated surrogate that
advances clocks by Manhattan distances and single-scalar dock
availabilities is systematically optimistic: it ignores collision
detours, dock-clearing maneuvers, waiting-line approach paths, and the
upstream holds triggered when waiting lines saturate. The state
features of Section~\ref{sec:features}---and therefore every decision
the policy makes---would then be conditioned on clocks that drift
further from reality as congestion grows, precisely in the regime where
scheduling decisions matter most. Rather than paying the computational
cost of embedding the full executor in the training loop, we close most
of this gap with two lightweight corrections fitted offline.

\subsection{Calibrated Travel and Dock-Wait Estimators}

\emph{Travel.} Realized leg times are modeled as an affine function of
the Manhattan distance, clamped from below by the metric itself:
\begin{equation}
\hat{t}(u,v)=\max\!\big(\Delta(u,v),\ \alpha\,\Delta(u,v)+\beta\big).
\label{eq:travelcal}
\end{equation}

\emph{Dock waiting.} The wait experienced by an AMR that becomes ready
at a dock $d$ at time $\rho$ is modeled as the committed availability
delay plus a penalty indexed by the queue depth at readiness,
\begin{equation}
\hat{w}_d(\rho)=\max\!\big(0,\ \mathrm{avail}_d-\rho\big)
+\varphi\!\big(\min(q_d(\rho),\,3)\big),
\label{eq:waitcal}
\end{equation}
where $q_d(\rho)=|\{e\in\mathcal{E}_d : c_e>\rho\}|$ counts committed
services unfinished at $\rho$, and
$\varphi:\{0,1,2,3^{+}\}\to\mathbb{R}_{\ge0}$ is a monotone lookup
table. The bounded table shape is deliberate: waiting-line detours and
upstream holds produce a cost cliff once the line saturates, a
nonlinearity that a linear congestion model cannot express.

\subsection{Fitting Procedure}

Calibration data are obtained by replaying complete schedules through
both simulators and aligning them per operation. Schedules are drawn
from heterogeneous sources---five dispatching-rule pairs and the
current learned policies---across instance sizes
$n\in\{10,20,40,60\}$, so that the fitted correction covers the state
distribution that policies actually visit. For every operation, the
raw surrogate supplies the Manhattan leg distance, the base wait, and
the queue depth at readiness; the executor's logged timeline is
partitioned, per AMR, into legs terminating at each service start,
within which movement entries sum to the realized travel and the
remainder is attributed to waiting. Executions containing infeasible
operations are discarded. The travel parameters $(\alpha,\beta)$ are
fitted by least squares and the penalties $\varphi(\cdot)$ as clipped
bin means with monotonicity enforced across depths.

In our environment, fitting on $10{,}360$ operation pairs from $159$
schedules yielded $\alpha=0.963$, $\beta=1.12$ ($R^{2}=0.91$) and
$\varphi=(0.04,\,0.04,\,1.66,\,3.57)$; the step between depths one and
two locates the waiting-line saturation cliff. The correction removes
the mean per-leg travel bias entirely ($+0.744\to0.000$ time units)
and halves the end-to-end makespan drift between surrogate and
executor (mean absolute error $24.3\to12.1$ on held-out schedules).
The calibration is a frozen artifact: it is fitted once before
training, adds no learned components to the training loop, and is
shared by \emph{all} policies and by the dispatching-rule baseline of
Section~\ref{sec:training}, so every method plans in the same
corrected world. After training converges, the fit can be re-estimated
from the new policies' schedules to confirm stability.

% =========================================================================

\section{Policy Optimization}
\label{sec:training}

\subsection{Objective and Gradient Estimator}

The policy $\pi_\theta$ is trained to minimize the expected executed
makespan $\mathbb{E}_{\pi_\theta}[C_{\max}(\Phi(\bm{a},\sigma))]$ with
the REINFORCE estimator \cite{williams1992reinforce}. During training,
episodes are generated by sampling from the masked policy
\eqref{eq:policy} on the calibrated surrogate; at validation and
deployment, actions are selected greedily. The gradient quality of
REINFORCE hinges on its baseline; we use a counterfactual baseline
built from a strong dispatching rule and evaluated by the executor
itself.

\subsection{Counterfactual Dispatch-Rule Baseline}
\label{sec:baseline}

Let $R$ denote a deterministic dispatching heuristic (job rule and AMR
rule; we use earliest-completion for both) and let
$R(\sigma_{1:t})$ denote the complete solution obtained by executing
the partial sequence $\sigma_{1:t}$ and completing all remaining
operations with $R$. Define the shorthand
$C[\cdot]=C_{\max}(\Phi(\cdot))$. The advantage of the action taken at
step $t$ is the improvement of committing that action and then
following the rule, relative to following the rule immediately:
\begin{equation}
A_t \;=\; C\big[R(\sigma_{1:t-1})\big]\;-\;
C\big[R(\sigma_{1:t-1}\oplus a_t)\big].
\label{eq:advantage}
\end{equation}
Both completions are evaluated by the collision-aware executor, so the
learning signal is grounded in true solution quality rather than in
surrogate clocks. Unlike episode-level rollout baselines
\cite{rennie2017self,kool2019attention}, \eqref{eq:advantage} assigns
credit per decision: an early pickup that congests a dock is penalized
at the step where it is committed, not diluted across the episode.
The price is $\mathcal{O}(n)$ executor evaluations per episode, which
dominates training cost; we consider the sharper credit assignment a
favorable trade at the instance sizes studied here, and
Section~\ref{sec:training-details} lists a cheaper actor--critic
alternative.

\subsection{Loss Function}

Advantages are standardized to zero mean and unit variance across all
steps of a batch of $B$ episodes. The per-batch loss combines the
policy gradient with a load-balancing shaping term and an entropy
bonus:
\begin{equation}
\mathcal{L}(\theta)=
-\frac{1}{B}\sum_{b=1}^{B}\sum_{t=1}^{2n}
\Big[\log\pi_\theta(a_t^b|s_t^b)\,
\big(\hat{A}_t^{\,b}+\lambda_{\mathrm{lb}}A_t^{\mathrm{lb},b}\big)
+\lambda_{H}\,\mathcal{H}_t^{b}\Big],
\label{eq:loss}
\end{equation}
where $\hat{A}_t$ is the standardized advantage,
$\mathcal{H}_t$ is the entropy of the masked policy over feasible
actions, and the shaping term
$A^{\mathrm{lb}}_t=\min_{k'}c_{k'}-c_{k}$ (for pickups; zero for
deliveries) compares the assignment count $c_k$ of the chosen AMR with
the least-loaded AMR, discouraging degenerate fleet utilization early
in training \cite{ng1999shaping}. We use
$\lambda_{\mathrm{lb}}=0.1$ and $\lambda_{H}=0.01$.

\begin{algorithm}[t]
\caption{Training with counterfactual dispatch baseline}
\label{alg:training}
\begin{algorithmic}[1]
\Require policy $\pi_\theta$, rule $R$, executor $\Phi$, surrogate $\hat{\Psi}$
\State fit calibration $(\alpha,\beta,\varphi)$ offline; freeze (Sec.~\ref{sec:calibration})
\For{$\mathrm{epoch}=1,\dots,E$}
  \For{$b=1,\dots,B$}
    \State sample instance; roll out $\sigma^b\!\sim\!\pi_\theta$ on $\hat{\Psi}$
    \For{$t=1,\dots,2n$}
      \State $A^b_t \gets C[R(\sigma^b_{1:t-1})]-C[R(\sigma^b_{1:t-1}\oplus a^b_t)]$
    \EndFor
  \EndFor
  \State standardize $\{A^b_t\}$ over the batch
  \State update $\theta$ by Adam on $\mathcal{L}(\theta)$ in \eqref{eq:loss}; clip $\lVert\nabla\rVert\le1$
  \If{$\mathrm{epoch}\bmod T_{\mathrm{val}}=0$}
    \State greedy-decode held-out instances; score by
    $C_{\max}+\kappa\cdot(\text{infeasible ops})$; keep best
  \EndIf
\EndFor
\end{algorithmic}
\end{algorithm}

\subsection{Training Protocol}
\label{sec:training-details}

Algorithm~\ref{alg:training} summarizes the procedure. Only the actor
parameters (encoders, interaction block, GIN layers, operation
embedding, and actor head) receive gradients under REINFORCE; the
critic \eqref{eq:critic} is reserved for an optional
proximal-policy-optimization variant \cite{schulman2017ppo} that
replaces \eqref{eq:advantage} with critic-based advantages when the
executor evaluations of the stepwise baseline are prohibitive. We
train with Adam \cite{kingma2015adam} (learning rate $3\times10^{-4}$,
batch $B=16$, gradient-norm clip $1.0$) for $2{,}000$ epochs. Every
$T_{\mathrm{val}}=50$ epochs the policy is evaluated by greedy decoding on a fixed
held-out validation set, disjoint from the training pool, and scored
by the executed makespan plus a penalty of $\kappa=10^{3}$ per
infeasible operation; the checkpoint with the best validation score is
retained. Two safeguards proved important in practice: the
identity-at-initialization residual branches of
Section~\ref{sec:interaction}, without which policy-gradient training
failed to converge, and computing the entropy in \eqref{eq:loss} over
the feasible-action support only, which keeps the bonus comparable as
the mask tightens toward the end of an episode.

% =========================================================================
% Suggested BibTeX entries (replace keys above with your bibliography):
%
% @inproceedings{vaswani2017attention, Vaswani et al., "Attention Is All
%   You Need", NeurIPS 2017}
% @inproceedings{xu2019gin, Xu, Hu, Leskovec, Jegelka, "How Powerful Are
%   Graph Neural Networks?", ICLR 2019}
% @inproceedings{xiong2020preln, Xiong et al., "On Layer Normalization in
%   the Transformer Architecture", ICML 2020}
% @inproceedings{bachlechner2021rezero, Bachlechner et al., "ReZero Is
%   All You Need: Fast Convergence at Large Depth", UAI 2021}
% @article{goyal2017accurate, Goyal et al., "Accurate, Large Minibatch
%   SGD: Training ImageNet in 1 Hour", arXiv:1706.02677}
% @inproceedings{zhang2020l2d, Zhang et al., "Learning to Dispatch for
%   Job Shop Scheduling via Deep Reinforcement Learning", NeurIPS 2020}
% @article{song2023fjsp, Song et al., "Flexible Job-Shop Scheduling via
%   Graph Neural Network and Deep Reinforcement Learning", IEEE TII 2023}
% @inproceedings{stern2019mapf, Stern et al., "Multi-Agent Pathfinding:
%   Definitions, Variants, and Benchmarks", SoCS 2019}
% @article{williams1992reinforce, Williams, "Simple Statistical
%   Gradient-Following Algorithms for Connectionist Reinforcement
%   Learning", Machine Learning 8, 1992}
% @inproceedings{rennie2017self, Rennie et al., "Self-Critical Sequence
%   Training for Image Captioning", CVPR 2017}
% @inproceedings{kool2019attention, Kool, van Hoof, Welling, "Attention,
%   Learn to Solve Routing Problems!", ICLR 2019}
% @inproceedings{ng1999shaping, Ng, Harada, Russell, "Policy Invariance
%   Under Reward Transformations: Theory and Application to Reward
%   Shaping", ICML 1999}
% @article{schulman2017ppo, Schulman et al., "Proximal Policy
%   Optimization Algorithms", arXiv:1707.06347}
% @inproceedings{kingma2015adam, Kingma, Ba, "Adam: A Method for
%   Stochastic Optimization", ICLR 2015}
% =========================================================================
 inside an IEEEtran document.
% Requires in the preamble:
%   \usepackage{amsmath,amssymb,bm}
%   \usepackage{booktabs,array,multirow,graphicx}
%   \usepackage{algorithm}
%   \usepackage{algpseudocode}
% Citation keys are placeholders -- suggested BibTeX entries are listed in
% comments at the end of this file. Replace \cite{...} keys with your own.
% "DockGNN" is a placeholder model name; rename globally if you prefer.
% Contains four sections: Problem Formulation, Architecture, Surrogate
% Calibration (sec:calibration), and Policy Optimization (sec:training).
% =========================================================================

\section{Problem Formulation}
\label{sec:problem}

\begin{table*}[t]
\caption{Summary of Notation}
\label{tab:notation}
\centering
\footnotesize
\begin{tabular}{@{}l p{0.295\textwidth} @{\qquad} l p{0.295\textwidth}@{}}
\toprule
Symbol & Meaning & Symbol & Meaning \\
\midrule
\multicolumn{2}{@{}l}{\itshape Problem setting} &
\multicolumn{2}{@{}l}{\itshape Architecture} \\
$\mathcal{W}$ & grid workspace (four-connected unit cells) &
  $\bm{x}^{A}_k,\bm{x}^{J}_j,\bm{x}^{D}_d$ & AMR / job / dock input features ($\mathbb{R}^{12},\mathbb{R}^{11},\mathbb{R}^{9}$) \\
$\Delta(u,v)$ & Manhattan distance between positions $u,v$ &
  $\bm{x}^{\mathrm{act}}_{\omega,k,j}$ & per-action decision features ($\mathbb{R}^{3}$) \\
$\mathcal{J},\ n$ & job set; number of jobs &
  $h$ & hidden width of all embeddings ($h{=}128$) \\
$\mathcal{M},\ m$ & AMR fleet; fleet size &
  $\bm{h}^{A}_k,\ \bm{h}^{D}_d$ & AMR and dock embeddings \\
$\mathcal{T}$ & set of material types &
  $\bar{\bm{h}}^{J}_j$ & job embedding before dock fusion \\
$\mathcal{D}^{\mathrm{in}},\mathcal{D}^{\mathrm{out}},\mathcal{D}$ & inbound docks; outbound docks; their union &
  $\bm{H}^{A},\bm{H}^{D}$ & stacked AMR / dock embedding matrices \\
$\tau_j$ & material type of job $j$ &
  $\mathrm{MHA},\mathrm{LN}$ & multi-head attention; layer normalization \\
$o_j,\ d_j$ & inbound and outbound dock of job $j$ &
  $\bm{g}^{(l)}_j$ & job embedding after fusion and $l$ GIN layers \\
$r_j$ & release time of job $j$ &
  $\mathcal{G}_t=(\mathcal{J},E_t)$ & contention--precedence graph at step $t$ \\
$p_j=p(\tau_j)$ & dock service duration of job $j$ &
  $\mathcal{N}(j),\ \deg_j$ & graph neighbors and degree of job $j$ \\
$\pi_j,\ \delta_j$ & pickup and delivery operations of job $j$ &
  $\epsilon^{(l)},\ L$ & learnable GIN self-weight; number of GIN layers \\
$Q$ & per-type AMR carrying capacity &
  $\bm{e}_{\omega}$ & learned operation-type embedding \\
$\bm{a},\ a_j$ & AMR assignment vector; AMR assigned to job $j$ &
  $d(\omega,j)$ & dock served by the action ($o_j$ if pickup, else $d_j$) \\
$\sigma,\ \sigma_t$ & operation sequence; prefix committed up to step $t$ &
  $z_{\omega,k,j}$ & logit of action $(\omega,k,j)$ \\
$s^{\pi}_j,c^{\pi}_j$ & start / completion of pickup service of job $j$ &
  $V_{\phi}$ & critic (state-value) network, parameters $\phi$ \\
$s^{\delta}_j,c^{\delta}_j$ & start / completion of delivery service of job $j$ &
  & \\
$t_k(u,v)$ & realized travel time of AMR $k$ from $u$ to $v$ &
  \multicolumn{2}{@{}l}{\itshape Training} \\
$\Phi$ & collision-aware executor &
  $R,\ R(\sigma_{1:t})$ & dispatching rule; completion of prefix $\sigma_{1:t}$ by $R$ \\
$C_{\max},\ F_k$ & makespan; finish time of AMR $k$ incl.\ depot return &
  $C[\cdot]$ & shorthand for $C_{\max}(\Phi(\cdot))$ \\
$W_d$ & number of physical waiting cells at dock $d$ &
  $A_t,\ \hat{A}_t$ & stepwise counterfactual advantage; standardized \\
\multicolumn{2}{@{}l}{\itshape Sequential decision process} &
  $A^{\mathrm{lb}}_t$ & load-balance shaping advantage \\
$s_t,\ a_t$ & state and action at decision epoch $t$ &
  $\lambda_{\mathrm{lb}},\ \lambda_{H}$ & shaping and entropy coefficients \\
$\omega$ & operation type, $\omega\in\{\textsc{pick},\textsc{deliver}\}$ &
  $\mathcal{H}_t$ & entropy of the masked policy at step $t$ \\
$\Lambda(s_t)$ & feasible-action set (mask) in state $s_t$ &
  $\mathcal{L}(\theta)$ & training loss \\
$\hat{\Psi}$ & calibrated surrogate transition model &
  $B,\ E$ & batch size (episodes); number of training epochs \\
$\mathcal{E}_d$ & committed service events $(s_e,c_e,j_e,\omega_e)$ at dock $d$ &
  $T_{\mathrm{val}},\ \kappa$ & validation interval; infeasibility penalty \\
$\hat{b}_k,\ x_k$ & next-available time and position of AMR $k$ &
  $\mathbf{1}[\cdot]$ & indicator function \\
$\rho$ & ready time of an AMR at a dock &
  & \\
$\hat{t}(u,v)$ & calibrated travel-time estimate \eqref{eq:travelcal} &
  & \\
$\hat{w}_d(\rho)$ & queue-aware dock-wait estimate \eqref{eq:waitcal} &
  & \\
$q_d(\rho)$ & services still unfinished at time $\rho$ at dock $d$ &
  & \\
$\alpha,\ \beta$ & travel-calibration scale and offset &
  & \\
$\varphi(\cdot)$ & dock-wait penalty table over queue depths &
  & \\
$\pi_{\theta}$ & scheduling policy with parameters $\theta$ &
  & \\
\bottomrule
\end{tabular}
\end{table*}

\subsection{System Description}

We consider the internal transfer problem of a cross-docking terminal
served by a fleet of autonomous mobile robots (AMRs). Table~\ref{tab:notation}
summarizes the notation used throughout the paper. The terminal floor
is modeled as a four-connected unit grid $\mathcal{W}$ in which a set of
\emph{inbound docks} $\mathcal{D}^{\mathrm{in}}$ is located along one
boundary and a set of \emph{outbound docks} (processing stations)
$\mathcal{D}^{\mathrm{out}}$ along the opposite boundary; the AMR depots
lie between them. Consistent with cross-docking practice, the floor is
free of static obstacles, so congestion arises from vehicle interference
and dock contention rather than from the layout itself. AMRs travel at
unit speed between adjacent cells, and the distance between two
positions $u,v$ is the Manhattan metric $\Delta(u,v)$.

A set of transfer jobs $\mathcal{J}=\{1,\dots,n\}$ must be moved across
the terminal by a fleet $\mathcal{M}=\{1,\dots,m\}$ of homogeneous AMRs.
Each job $j\in\mathcal{J}$ represents one unit load of a material type
$\tau_j\in\mathcal{T}$ that becomes available at its inbound dock
$o_j\in\mathcal{D}^{\mathrm{in}}$ at release time $r_j\ge 0$ and must be
delivered to a designated outbound dock
$d_j\in\mathcal{D}^{\mathrm{out}}$. Handling a unit load of type $\tau$
occupies a dock for a type-dependent service duration $p(\tau)$; we
write $p_j=p(\tau_j)$ and assume the loading service at the inbound dock
and the unloading/processing service at the outbound dock take the same
time, reflecting symmetric handling of a unit load. Each job therefore
consists of an ordered pair of operations: a \emph{pickup} operation
$\pi_j$ executed at $o_j$ and a \emph{delivery} operation $\delta_j$
executed at $d_j$, both performed by the same AMR (no transshipment).
An AMR may consolidate loads subject to a per-type carrying capacity
$Q$: at no time may it carry more than $Q$ units of any single material
type.

Docks are exclusive single-server resources: at most one AMR may be in
service at a dock at any time, and services are non-preemptive. An AMR
arriving at an occupied dock queues in a bounded set of waiting cells
adjacent to that dock; when the waiting line is saturated, the AMR must
hold at an upstream position and re-approach later. Finally, AMR motion
must be collision-free: at any unit time step, each grid cell may be
occupied by at most one AMR.

\subsection{Formal Statement}

A solution is a pair $(\bm{a},\sigma)$, where
$\bm{a}=(a_1,\dots,a_n)\in\mathcal{M}^{n}$ assigns each job to an AMR
and $\sigma$ is a total order over the $2n$ operations
$\{\pi_j,\delta_j\}_{j\in\mathcal{J}}$ that determines the sequence in
which operations are committed. Let $s^{\pi}_j,c^{\pi}_j$
(resp.\ $s^{\delta}_j,c^{\delta}_j$) denote the service start and
completion times of $\pi_j$ (resp.\ $\delta_j$) induced by executing
$(\bm{a},\sigma)$. A feasible execution satisfies:
\begin{align}
c^{\pi}_j &= s^{\pi}_j + p_j, \qquad
c^{\delta}_j = s^{\delta}_j + p_j,
&& \forall j\in\mathcal{J}, \label{eq:nonpreempt}\\
s^{\pi}_j &\ge r_j,
&& \forall j\in\mathcal{J}, \label{eq:release}\\
s^{\delta}_j &\ge c^{\pi}_j,
&& \forall j\in\mathcal{J}, \label{eq:precedence}
\end{align}
\begin{align}
\big|\{\,j : a_j{=}k,\ \tau_j{=}\tau,\ c^{\pi}_j \le u < c^{\delta}_j \}\big| &\le Q,
&& \forall k,\tau,u, \label{eq:capacity}\\
\big[s_e, c_e\big) \cap \big[s_{e'}, c_{e'}\big) &= \emptyset,
&& \forall e\ne e' \ \text{at the same dock}, \label{eq:exclusive}
\end{align}
where \eqref{eq:nonpreempt} enforces non-preemptive services,
\eqref{eq:release} release dates, \eqref{eq:precedence} pickup-before-%
delivery precedence, \eqref{eq:capacity} the per-type carrying capacity,
and \eqref{eq:exclusive} dock exclusivity over all service events $e$.
In addition, for any two operations $u\prec v$ executed consecutively by
the same AMR $k$,
\begin{equation}
s_v \;\ge\; c_u + t_k\big(\mathrm{loc}(u),\mathrm{loc}(v)\big),
\label{eq:travel}
\end{equation}
where $t_k(\cdot,\cdot)\ge\Delta(\cdot,\cdot)$ is the realized travel
time of AMR $k$, which may exceed the Manhattan distance due to
collision avoidance, queuing detours, and dock-clearing maneuvers.

The realized travel and waiting times are not fixed constants but are
determined by a high-fidelity \emph{executor} $\Phi$, which routes every
AMR leg with space--time $\mathrm{A}^{*}$ search over a shared
reservation table (at most one AMR per cell per time step), manages the
bounded dock waiting lines, and inserts upstream holding maneuvers when
lines are saturated. The problem is therefore a simulation-based
optimization: find
\begin{equation}
(\bm{a}^{*},\sigma^{*}) \;=\;
\arg\min_{(\bm{a},\sigma)} \; C_{\max}\big(\Phi(\bm{a},\sigma)\big),
\label{eq:objective}
\end{equation}
where the makespan $C_{\max}=\max_{k\in\mathcal{M}} F_k$ and $F_k$ is
the time at which AMR $k$ completes its final activity, including its
return trip to the depot. The problem generalizes both the capacitated
pickup-and-delivery problem and parallel-machine scheduling with
sequence-dependent setups, and is therefore NP-hard; the coupling to
collision-free multi-agent routing further compounds its difficulty
\cite{stern2019mapf}.

\subsection{Sequential Decision Formulation}
\label{sec:mdp}

We cast the construction of $(\bm{a},\sigma)$ as an episodic sequential
decision process with exactly $2n$ decision epochs, one per committed
operation. At epoch $t$, the state
\begin{equation}
s_t=\big(\{\bm{x}^{A}_{k}\},\, \{\bm{x}^{J}_{j}\},\, \{\bm{x}^{D}_{d}\},\,
\{\mathcal{E}_d\},\, \sigma_{t}\big)
\label{eq:state}
\end{equation}
collects the AMR states (position, next-available time, per-type
inventory), the job states (unpicked / onboard / delivered, carrier
identity), the dock states, the set of committed dock service events
$\mathcal{E}_d=\{(s_e,c_e,j_e,\omega_e)\}$ at each dock $d$, and the
partial operation sequence $\sigma_t$.

An action is a triple
\begin{equation}
a_t=(\omega,k,j)\in\{\textsc{pick},\textsc{deliver}\}\times\mathcal{M}\times\mathcal{J},
\label{eq:action}
\end{equation}
giving a flat action space of size $2mn$ restricted by a feasibility
mask $\Lambda(s_t)$: $(\textsc{pick},k,j)$ is legal iff $j$ is unpicked
and AMR $k$ carries fewer than $Q$ units of type $\tau_j$;
$(\textsc{deliver},k,j)$ is legal iff $j$ is onboard AMR $k$. The mask
guarantees by construction that every complete episode yields a feasible
solution, so no repair mechanism is required.

Transitions follow a deterministic surrogate model $\hat{\Psi}$ that
advances the acting AMR's clock using a \emph{calibrated} travel-time
estimate $\hat{t}(\cdot,\cdot)$ and a queue-aware dock-wait estimate
$\hat{w}_d(\cdot)$ fitted offline against the executor $\Phi$
(Section~\ref{sec:calibration}); committing an action appends the
corresponding service window to $\mathcal{E}_d$. For a pickup, the
service start of job $j$ by AMR $k$ located at $x_k$ with next-available
time $\hat{b}_k$ is
\begin{equation}
s^{\pi}_j=\underbrace{\max\!\big(\hat{b}_k+\hat{t}(x_k,o_j),\, r_j\big)}_{\text{ready time }\rho}
\;+\;\hat{w}_{o_j}(\rho),
\label{eq:surrogate}
\end{equation}
and analogously for deliveries at $d_j$. After $2n$ steps the completed
solution is evaluated by the executor, and the learning objective is to
minimize the expected executed makespan
$\mathbb{E}_{\pi_\theta}\!\left[C_{\max}(\Phi(\bm{a},\sigma))\right]$;
the policy-gradient training procedure and its dispatch-rule baseline
are described in Section~\ref{sec:training}.

% =========================================================================

\section{Dock-Aware Graph Scheduler}
\label{sec:architecture}

\subsection{Overview and Design Rationale}

Existing learning-to-dispatch architectures for shop scheduling encode
jobs (operations) as graph nodes and machines as flat feature vectors
\cite{zhang2020l2d,song2023fjsp}. Applied to cross-docking, this design
has three blind spots. First, the binding resources are the docks: with
$m$ AMRs contending for $|\mathcal{D}^{\mathrm{in}}|+|\mathcal{D}^{\mathrm{out}}|$
single-server docks, queuing at docks---not processing---dominates the
makespan, yet job-only graphs represent dock congestion at best as
scalar features replicated across jobs. Second, the temporal structure
of dock occupancy (which service windows are already committed, and
when) is invisible. Third, the AMR fleet is encoded independently
per vehicle, so mutual interference cannot be anticipated.

The proposed scheduler, hereafter \textsc{DockGNN}, addresses all three
by treating docks as first-class graph entities. It comprises
(i)~heterogeneous encoders for AMR, job, and dock nodes;
(ii)~an AMR--dock interaction block that exchanges information between
the fleet and the dock set;
(iii)~a job \emph{contention--precedence} graph processed by
degree-normalized graph isomorphism network (GIN) layers
\cite{xu2019gin}; and
(iv)~a factored actor head that scores each masked action from the
operation type, the acting AMR, the target job, the operation-specific
dock, and a small set of per-action decision features. The full model
has $6.2\times10^{5}$ parameters ($5.7\times10^{5}$ actor,
$4.9\times10^{4}$ critic) with hidden width $h=128$ throughout.

% TODO: architecture overview figure
% \begin{figure*}[t]
%   \centering
%   \includegraphics[width=\textwidth]{figures/dockgnn_architecture}
%   \caption{Overview of \textsc{DockGNN}. Dock, AMR, and job features are
%   encoded separately; an AMR--dock attention block exchanges congestion
%   information; jobs fuse their inbound/outbound dock embeddings and pass
%   messages over the contention--precedence graph; the actor scores each
%   feasible (operation, AMR, job) triple.}
%   \label{fig:architecture}
% \end{figure*}

\subsection{Heterogeneous State Representation}
\label{sec:features}

At every decision epoch the state \eqref{eq:state} is rendered into four
feature families, summarized in Table~\ref{tab:features}. All features
are normalized to comparable magnitudes by fixed scales.

\emph{Dock nodes.} Each dock $d$ maintains its committed service events
$\mathcal{E}_d$. Writing $t$ for the current decision time, the derived
features are the availability delay
$\max(0,\mathrm{avail}_d-t)$, the number of queued future services
$q_d=|\{e\in\mathcal{E}_d : s_e>t\}|$, the waiting-line saturation
$\min(q_d/W_d,1)$ with $W_d$ the number of physical waiting cells, the
residual time of the service in progress, and the committed workload
$\sum_{e\,\text{unfinished}}\big(c_e-\max(s_e,t)\big)$, together with
the dock class (inbound/outbound) and its coordinates.

\emph{AMR nodes.} Each AMR encodes its busy flag and time-to-free,
per-type inventory ratios and residual capacity, the number of loads
currently onboard, spatial context (coordinates, distance to the
nearest dock, a near-dock indicator), and a local fleet density---the
fraction of other AMRs within a two-cell radius---which is the model's
proxy for imminent interference.

\emph{Job nodes.} Each job encodes its service duration, waiting time
since release, lifecycle status (unpicked/onboard/delivered), carrier
identity, material type, dock indices, and a completion-time estimate
computed with the calibrated travel model
(Section~\ref{sec:calibration}).

\emph{Per-action features.} For every candidate triple $(\omega,k,j)$,
three decision-grade quantities are computed directly: the calibrated
travel time from AMR $k$'s position to the service dock, the calibrated
downstream distance ($o_j$ to $d_j$ for pickups; zero for deliveries),
and the projected queue-aware dock wait $\hat{w}(\cdot)$ of
\eqref{eq:surrogate}. Providing these estimates explicitly relieves the
network from recovering travel arithmetic from coordinates through
several nonlinear layers.

\begin{table}[t]
\caption{Node and Action Features of \textsc{DockGNN}}
\label{tab:features}
\centering
\footnotesize
\begin{tabular}{@{}llc@{}}
\toprule
& Feature & Scale \\
\midrule
\multirow{9}{*}{\rotatebox{90}{AMR $\bm{x}^{A}_k\!\in\!\mathbb{R}^{12}$}}
& busy indicator & -- \\
& time until available & $1/100$ \\
& inventory ratio, per type ($\times3$) & $1/Q$ \\
& residual capacity ratio & $1/(|\mathcal{T}|Q)$ \\
& loads onboard & $1/10$ \\
& near-dock indicator & -- \\
& distance to nearest dock & $1/20$ \\
& local fleet density ($\le\!2$ cells) & $1/(m{-}1)$ \\
& position $(x,y)$ ($\times2$) & $1/10$ \\
\midrule
\multirow{9}{*}{\rotatebox{90}{Job $\bm{x}^{J}_j\!\in\!\mathbb{R}^{11}$}}
& existence flag & -- \\
& service duration $p_j$ & $1/25$ \\
& waiting time since release & $1/100$ \\
& carrier index (0 if unpicked) & $1/m$ \\
& lifecycle status $\{0,\tfrac12,1\}$ & -- \\
& material type one-hot ($\times3$) & -- \\
& inbound dock index & $1/|\mathcal{D}^{\mathrm{in}}|$ \\
& outbound dock index & $1/|\mathcal{D}^{\mathrm{out}}|$ \\
& completion-time estimate & $1/500$ \\
\midrule
\multirow{7}{*}{\rotatebox{90}{Dock $\bm{x}^{D}_d\!\in\!\mathbb{R}^{9}$}}
& inbound / outbound class ($\times2$) & -- \\
& availability delay & $1/100$ \\
& queued services $q_d$ & $1/m$ \\
& waiting-line saturation & -- \\
& residual service in progress & $1/100$ \\
& committed workload & $1/500$ \\
& position $(x,y)$ ($\times2$) & $1/10$ \\
\midrule
\multirow{3}{*}{\rotatebox{90}{Act.\ $\mathbb{R}^{3}$}}
& calibrated travel to service dock & $1/20$ \\
& calibrated downstream distance & $1/20$ \\
& projected queue-aware dock wait & $1/100$ \\
\bottomrule
\end{tabular}
\end{table}

\subsection{Entity Encoders and AMR--Dock Interaction}
\label{sec:interaction}

AMR and dock features pass through two-layer perceptrons and job
features through a linear map,
\begin{equation}
\bm{h}^{A}_k=\mathrm{MLP}_{A}(\bm{x}^{A}_k),\quad
\bm{h}^{D}_d=\mathrm{MLP}_{D}(\bm{x}^{D}_d),\quad
\bar{\bm{h}}^{J}_j=\bm{W}_{\!J}\,\bm{x}^{J}_j,
\end{equation}
all in $\mathbb{R}^{h}$. Fleet--dock coupling is then established by a
single pre-normalization residual attention round
\cite{vaswani2017attention,xiong2020preln} consisting of three
sublayers: docks attend to the fleet, the fleet attends to itself, and
the fleet attends to the updated docks,
\begin{align}
\bm{H}^{D} &\leftarrow \bm{H}^{D} +
\mathrm{MHA}\big(\mathrm{LN}(\bm{H}^{D}),\,\mathrm{LN}(\bm{H}^{A}),\,\mathrm{LN}(\bm{H}^{A})\big),
\label{eq:dock-from-amr}\\
\bm{H}^{A} &\leftarrow \bm{H}^{A} +
\mathrm{MHA}\big(\mathrm{LN}(\bm{H}^{A}),\,\mathrm{LN}(\bm{H}^{A}),\,\mathrm{LN}(\bm{H}^{A})\big),
\label{eq:amr-self}\\
\bm{H}^{A} &\leftarrow \bm{H}^{A} +
\mathrm{MHA}\big(\mathrm{LN}(\bm{H}^{A}),\,\mathrm{LN}(\bm{H}^{D}),\,\mathrm{LN}(\bm{H}^{D})\big),
\label{eq:amr-from-dock}
\end{align}
with four heads per attention. Through
\eqref{eq:dock-from-amr}--\eqref{eq:amr-from-dock}, a dock's embedding
reflects which vehicles are loaded and heading toward it \emph{before}
jobs consume that embedding, and each AMR perceives both the remaining
fleet and network-wide dock congestion.

\emph{Stability.} The output projections of all three attention
sublayers are zero-initialized, making the entire block an exact
identity at the start of training in the spirit of residual-branch
zero initialization \cite{bachlechner2021rezero,goyal2017accurate}.
This is not cosmetic: the embedding stream entering the block has small
magnitude relative to the unit-scale outputs that randomly initialized
attention produces from layer-normalized inputs, and without
zero initialization the three stacked residual additions amplify the
fleet embeddings by more than an order of magnitude at initialization,
which we found to prevent policy-gradient training from converging.

\subsection{Job Contention--Precedence Graph}
\label{sec:graph}

Jobs interact through two mechanisms with different lifetimes:
\emph{future contention} for shared docks, and \emph{committed
precedence} induced by scheduling decisions already made. Both are
encoded in a single dynamic graph $\mathcal{G}_t=(\mathcal{J},E_t)$
with symmetric edges over the job nodes:
(i)~a \emph{contention edge} joins every pair of unfinished jobs that
share an inbound or an outbound dock, so resource competition is
visible from the first decision onward; and
(ii)~\emph{order edges} join consecutively scheduled jobs on the same
AMR and on the same outbound dock, incrementally densifying the graph
as $\sigma_t$ grows.

Before message passing, each job fuses the embeddings of its two docks,
\begin{equation}
\bm{g}^{(0)}_j=\mathrm{MLP}_{F}\!\big(\big[\,\bar{\bm{h}}^{J}_j \,\|\,
\bm{h}^{D}_{o_j} \,\|\, \bm{h}^{D}_{d_j}\big]\big),
\label{eq:fusion}
\end{equation}
so every job representation is conditioned on the live congestion state
of exactly the two resources it will occupy. Message passing then
applies $L=3$ residual GIN layers with \emph{degree-normalized}
aggregation and layer normalization,
\begin{equation}
\bm{g}^{(l+1)}_j=\bm{g}^{(l)}_j+
\mathrm{LN}\!\left(\mathrm{MLP}^{(l)}\!\Big(
(1{+}\epsilon^{(l)})\,\bm{g}^{(l)}_j+
\tfrac{1}{\deg_j}\!\!\sum_{i\in\mathcal{N}(j)}\!\bm{g}^{(l)}_i
\Big)\right)\!.
\label{eq:gin}
\end{equation}
Mean aggregation deliberately trades the injectivity of sum-based GIN
\cite{xu2019gin} for scale invariance: contention edges form dock
cliques whose size grows linearly with $n/|\mathcal{D}|$, and
sum aggregation would make message magnitudes---and hence the
activation statistics of every downstream layer---depend on the
instance size, undermining transfer across job counts. With mean
aggregation the same trained network processes $n{=}10$ and $n{=}100$
instances at identical activation scales.

\subsection{Action Scoring and Value Estimation}
\label{sec:actor}

Let $\bm{e}_{\omega}\in\mathbb{R}^{h}$ be a learned embedding of the
operation type and let $d(\omega,j)=o_j$ if $\omega=\textsc{pick}$ and
$d_j$ otherwise denote the dock at which the candidate action would be
served. The actor scores every triple by
\begin{equation}
z_{\omega,k,j}=\mathrm{MLP}_{\pi}\!\big(\big[\,\bm{e}_{\omega}\,\|\,
\bm{h}^{A}_{k}\,\|\,\bm{g}^{(L)}_{j}\,\|\,
\bm{h}^{D}_{d(\omega,j)}\,\|\,\bm{x}^{\mathrm{act}}_{\omega,k,j}\big]\big),
\label{eq:logit}
\end{equation}
and the policy is the masked softmax
\begin{equation}
\pi_{\theta}(a\,|\,s_t)=
\frac{\exp(z_{a})\,\mathbf{1}[a\in\Lambda(s_t)]}
{\sum_{a'\in\Lambda(s_t)}\exp(z_{a'})}.
\label{eq:policy}
\end{equation}
The critic estimates the state value from pooled entity summaries,
\begin{equation}
V_{\phi}(s_t)=\mathrm{MLP}_{V}\!\big(\big[\,
\overline{\bm{g}}^{\mathrm{(unf)}} \,\|\, \overline{\bm{h}}^{A} \,\|\,
\overline{\bm{h}}^{D}\big]\big),
\label{eq:critic}
\end{equation}
where $\overline{\bm{g}}^{\mathrm{(unf)}}$ averages the embeddings of
unfinished jobs only. A full episode requires exactly $2n$ forward
passes; each pass scores $2mn$ actions in a single batched evaluation
of \eqref{eq:logit}, and all components are permutation-invariant in
the job and AMR sets, so the architecture applies unchanged to any
instance size.

% =========================================================================

\section{Surrogate Model Calibration}
\label{sec:calibration}

\subsection{The Two-Simulator Gap}

The policy constructs schedules on the surrogate $\hat{\Psi}$ of
Section~\ref{sec:mdp}, whereas solution quality is defined by the
executor $\Phi$ of \eqref{eq:objective}. An uncalibrated surrogate that
advances clocks by Manhattan distances and single-scalar dock
availabilities is systematically optimistic: it ignores collision
detours, dock-clearing maneuvers, waiting-line approach paths, and the
upstream holds triggered when waiting lines saturate. The state
features of Section~\ref{sec:features}---and therefore every decision
the policy makes---would then be conditioned on clocks that drift
further from reality as congestion grows, precisely in the regime where
scheduling decisions matter most. Rather than paying the computational
cost of embedding the full executor in the training loop, we close most
of this gap with two lightweight corrections fitted offline.

\subsection{Calibrated Travel and Dock-Wait Estimators}

\emph{Travel.} Realized leg times are modeled as an affine function of
the Manhattan distance, clamped from below by the metric itself:
\begin{equation}
\hat{t}(u,v)=\max\!\big(\Delta(u,v),\ \alpha\,\Delta(u,v)+\beta\big).
\label{eq:travelcal}
\end{equation}

\emph{Dock waiting.} The wait experienced by an AMR that becomes ready
at a dock $d$ at time $\rho$ is modeled as the committed availability
delay plus a penalty indexed by the queue depth at readiness,
\begin{equation}
\hat{w}_d(\rho)=\max\!\big(0,\ \mathrm{avail}_d-\rho\big)
+\varphi\!\big(\min(q_d(\rho),\,3)\big),
\label{eq:waitcal}
\end{equation}
where $q_d(\rho)=|\{e\in\mathcal{E}_d : c_e>\rho\}|$ counts committed
services unfinished at $\rho$, and
$\varphi:\{0,1,2,3^{+}\}\to\mathbb{R}_{\ge0}$ is a monotone lookup
table. The bounded table shape is deliberate: waiting-line detours and
upstream holds produce a cost cliff once the line saturates, a
nonlinearity that a linear congestion model cannot express.

\subsection{Fitting Procedure}

Calibration data are obtained by replaying complete schedules through
both simulators and aligning them per operation. Schedules are drawn
from heterogeneous sources---five dispatching-rule pairs and the
current learned policies---across instance sizes
$n\in\{10,20,40,60\}$, so that the fitted correction covers the state
distribution that policies actually visit. For every operation, the
raw surrogate supplies the Manhattan leg distance, the base wait, and
the queue depth at readiness; the executor's logged timeline is
partitioned, per AMR, into legs terminating at each service start,
within which movement entries sum to the realized travel and the
remainder is attributed to waiting. Executions containing infeasible
operations are discarded. The travel parameters $(\alpha,\beta)$ are
fitted by least squares and the penalties $\varphi(\cdot)$ as clipped
bin means with monotonicity enforced across depths.

In our environment, fitting on $10{,}360$ operation pairs from $159$
schedules yielded $\alpha=0.963$, $\beta=1.12$ ($R^{2}=0.91$) and
$\varphi=(0.04,\,0.04,\,1.66,\,3.57)$; the step between depths one and
two locates the waiting-line saturation cliff. The correction removes
the mean per-leg travel bias entirely ($+0.744\to0.000$ time units)
and halves the end-to-end makespan drift between surrogate and
executor (mean absolute error $24.3\to12.1$ on held-out schedules).
The calibration is a frozen artifact: it is fitted once before
training, adds no learned components to the training loop, and is
shared by \emph{all} policies and by the dispatching-rule baseline of
Section~\ref{sec:training}, so every method plans in the same
corrected world. After training converges, the fit can be re-estimated
from the new policies' schedules to confirm stability.

% =========================================================================

\section{Policy Optimization}
\label{sec:training}

\subsection{Objective and Gradient Estimator}

The policy $\pi_\theta$ is trained to minimize the expected executed
makespan $\mathbb{E}_{\pi_\theta}[C_{\max}(\Phi(\bm{a},\sigma))]$ with
the REINFORCE estimator \cite{williams1992reinforce}. During training,
episodes are generated by sampling from the masked policy
\eqref{eq:policy} on the calibrated surrogate; at validation and
deployment, actions are selected greedily. The gradient quality of
REINFORCE hinges on its baseline; we use a counterfactual baseline
built from a strong dispatching rule and evaluated by the executor
itself.

\subsection{Counterfactual Dispatch-Rule Baseline}
\label{sec:baseline}

Let $R$ denote a deterministic dispatching heuristic (job rule and AMR
rule; we use earliest-completion for both) and let
$R(\sigma_{1:t})$ denote the complete solution obtained by executing
the partial sequence $\sigma_{1:t}$ and completing all remaining
operations with $R$. Define the shorthand
$C[\cdot]=C_{\max}(\Phi(\cdot))$. The advantage of the action taken at
step $t$ is the improvement of committing that action and then
following the rule, relative to following the rule immediately:
\begin{equation}
A_t \;=\; C\big[R(\sigma_{1:t-1})\big]\;-\;
C\big[R(\sigma_{1:t-1}\oplus a_t)\big].
\label{eq:advantage}
\end{equation}
Both completions are evaluated by the collision-aware executor, so the
learning signal is grounded in true solution quality rather than in
surrogate clocks. Unlike episode-level rollout baselines
\cite{rennie2017self,kool2019attention}, \eqref{eq:advantage} assigns
credit per decision: an early pickup that congests a dock is penalized
at the step where it is committed, not diluted across the episode.
The price is $\mathcal{O}(n)$ executor evaluations per episode, which
dominates training cost; we consider the sharper credit assignment a
favorable trade at the instance sizes studied here, and
Section~\ref{sec:training-details} lists a cheaper actor--critic
alternative.

\subsection{Loss Function}

Advantages are standardized to zero mean and unit variance across all
steps of a batch of $B$ episodes. The per-batch loss combines the
policy gradient with a load-balancing shaping term and an entropy
bonus:
\begin{equation}
\mathcal{L}(\theta)=
-\frac{1}{B}\sum_{b=1}^{B}\sum_{t=1}^{2n}
\Big[\log\pi_\theta(a_t^b|s_t^b)\,
\big(\hat{A}_t^{\,b}+\lambda_{\mathrm{lb}}A_t^{\mathrm{lb},b}\big)
+\lambda_{H}\,\mathcal{H}_t^{b}\Big],
\label{eq:loss}
\end{equation}
where $\hat{A}_t$ is the standardized advantage,
$\mathcal{H}_t$ is the entropy of the masked policy over feasible
actions, and the shaping term
$A^{\mathrm{lb}}_t=\min_{k'}c_{k'}-c_{k}$ (for pickups; zero for
deliveries) compares the assignment count $c_k$ of the chosen AMR with
the least-loaded AMR, discouraging degenerate fleet utilization early
in training \cite{ng1999shaping}. We use
$\lambda_{\mathrm{lb}}=0.1$ and $\lambda_{H}=0.01$.

\begin{algorithm}[t]
\caption{Training with counterfactual dispatch baseline}
\label{alg:training}
\begin{algorithmic}[1]
\Require policy $\pi_\theta$, rule $R$, executor $\Phi$, surrogate $\hat{\Psi}$
\State fit calibration $(\alpha,\beta,\varphi)$ offline; freeze (Sec.~\ref{sec:calibration})
\For{$\mathrm{epoch}=1,\dots,E$}
  \For{$b=1,\dots,B$}
    \State sample instance; roll out $\sigma^b\!\sim\!\pi_\theta$ on $\hat{\Psi}$
    \For{$t=1,\dots,2n$}
      \State $A^b_t \gets C[R(\sigma^b_{1:t-1})]-C[R(\sigma^b_{1:t-1}\oplus a^b_t)]$
    \EndFor
  \EndFor
  \State standardize $\{A^b_t\}$ over the batch
  \State update $\theta$ by Adam on $\mathcal{L}(\theta)$ in \eqref{eq:loss}; clip $\lVert\nabla\rVert\le1$
  \If{$\mathrm{epoch}\bmod T_{\mathrm{val}}=0$}
    \State greedy-decode held-out instances; score by
    $C_{\max}+\kappa\cdot(\text{infeasible ops})$; keep best
  \EndIf
\EndFor
\end{algorithmic}
\end{algorithm}

\subsection{Training Protocol}
\label{sec:training-details}

Algorithm~\ref{alg:training} summarizes the procedure. Only the actor
parameters (encoders, interaction block, GIN layers, operation
embedding, and actor head) receive gradients under REINFORCE; the
critic \eqref{eq:critic} is reserved for an optional
proximal-policy-optimization variant \cite{schulman2017ppo} that
replaces \eqref{eq:advantage} with critic-based advantages when the
executor evaluations of the stepwise baseline are prohibitive. We
train with Adam \cite{kingma2015adam} (learning rate $3\times10^{-4}$,
batch $B=16$, gradient-norm clip $1.0$) for $2{,}000$ epochs. Every
$T_{\mathrm{val}}=50$ epochs the policy is evaluated by greedy decoding on a fixed
held-out validation set, disjoint from the training pool, and scored
by the executed makespan plus a penalty of $\kappa=10^{3}$ per
infeasible operation; the checkpoint with the best validation score is
retained. Two safeguards proved important in practice: the
identity-at-initialization residual branches of
Section~\ref{sec:interaction}, without which policy-gradient training
failed to converge, and computing the entropy in \eqref{eq:loss} over
the feasible-action support only, which keeps the bonus comparable as
the mask tightens toward the end of an episode.

% =========================================================================
% Suggested BibTeX entries (replace keys above with your bibliography):
%
% @inproceedings{vaswani2017attention, Vaswani et al., "Attention Is All
%   You Need", NeurIPS 2017}
% @inproceedings{xu2019gin, Xu, Hu, Leskovec, Jegelka, "How Powerful Are
%   Graph Neural Networks?", ICLR 2019}
% @inproceedings{xiong2020preln, Xiong et al., "On Layer Normalization in
%   the Transformer Architecture", ICML 2020}
% @inproceedings{bachlechner2021rezero, Bachlechner et al., "ReZero Is
%   All You Need: Fast Convergence at Large Depth", UAI 2021}
% @article{goyal2017accurate, Goyal et al., "Accurate, Large Minibatch
%   SGD: Training ImageNet in 1 Hour", arXiv:1706.02677}
% @inproceedings{zhang2020l2d, Zhang et al., "Learning to Dispatch for
%   Job Shop Scheduling via Deep Reinforcement Learning", NeurIPS 2020}
% @article{song2023fjsp, Song et al., "Flexible Job-Shop Scheduling via
%   Graph Neural Network and Deep Reinforcement Learning", IEEE TII 2023}
% @inproceedings{stern2019mapf, Stern et al., "Multi-Agent Pathfinding:
%   Definitions, Variants, and Benchmarks", SoCS 2019}
% @article{williams1992reinforce, Williams, "Simple Statistical
%   Gradient-Following Algorithms for Connectionist Reinforcement
%   Learning", Machine Learning 8, 1992}
% @inproceedings{rennie2017self, Rennie et al., "Self-Critical Sequence
%   Training for Image Captioning", CVPR 2017}
% @inproceedings{kool2019attention, Kool, van Hoof, Welling, "Attention,
%   Learn to Solve Routing Problems!", ICLR 2019}
% @inproceedings{ng1999shaping, Ng, Harada, Russell, "Policy Invariance
%   Under Reward Transformations: Theory and Application to Reward
%   Shaping", ICML 1999}
% @article{schulman2017ppo, Schulman et al., "Proximal Policy
%   Optimization Algorithms", arXiv:1707.06347}
% @inproceedings{kingma2015adam, Kingma, Ba, "Adam: A Method for
%   Stochastic Optimization", ICLR 2015}
% =========================================================================
 inside an IEEEtran document.
% Requires in the preamble:
%   \usepackage{amsmath,amssymb,bm}
%   \usepackage{booktabs,array,multirow,graphicx}
%   \usepackage{algorithm}
%   \usepackage{algpseudocode}
% Citation keys are placeholders -- suggested BibTeX entries are listed in
% comments at the end of this file. Replace \cite{...} keys with your own.
% "DockGNN" is a placeholder model name; rename globally if you prefer.
% Contains four sections: Problem Formulation, Architecture, Surrogate
% Calibration (sec:calibration), and Policy Optimization (sec:training).
% =========================================================================

\section{Problem Formulation}
\label{sec:problem}

\begin{table*}[t]
\caption{Summary of Notation}
\label{tab:notation}
\centering
\footnotesize
\begin{tabular}{@{}l p{0.295\textwidth} @{\qquad} l p{0.295\textwidth}@{}}
\toprule
Symbol & Meaning & Symbol & Meaning \\
\midrule
\multicolumn{2}{@{}l}{\itshape Problem setting} &
\multicolumn{2}{@{}l}{\itshape Architecture} \\
$\mathcal{W}$ & grid workspace (four-connected unit cells) &
  $\bm{x}^{A}_k,\bm{x}^{J}_j,\bm{x}^{D}_d$ & AMR / job / dock input features ($\mathbb{R}^{12},\mathbb{R}^{11},\mathbb{R}^{9}$) \\
$\Delta(u,v)$ & Manhattan distance between positions $u,v$ &
  $\bm{x}^{\mathrm{act}}_{\omega,k,j}$ & per-action decision features ($\mathbb{R}^{3}$) \\
$\mathcal{J},\ n$ & job set; number of jobs &
  $h$ & hidden width of all embeddings ($h{=}128$) \\
$\mathcal{M},\ m$ & AMR fleet; fleet size &
  $\bm{h}^{A}_k,\ \bm{h}^{D}_d$ & AMR and dock embeddings \\
$\mathcal{T}$ & set of material types &
  $\bar{\bm{h}}^{J}_j$ & job embedding before dock fusion \\
$\mathcal{D}^{\mathrm{in}},\mathcal{D}^{\mathrm{out}},\mathcal{D}$ & inbound docks; outbound docks; their union &
  $\bm{H}^{A},\bm{H}^{D}$ & stacked AMR / dock embedding matrices \\
$\tau_j$ & material type of job $j$ &
  $\mathrm{MHA},\mathrm{LN}$ & multi-head attention; layer normalization \\
$o_j,\ d_j$ & inbound and outbound dock of job $j$ &
  $\bm{g}^{(l)}_j$ & job embedding after fusion and $l$ GIN layers \\
$r_j$ & release time of job $j$ &
  $\mathcal{G}_t=(\mathcal{J},E_t)$ & contention--precedence graph at step $t$ \\
$p_j=p(\tau_j)$ & dock service duration of job $j$ &
  $\mathcal{N}(j),\ \deg_j$ & graph neighbors and degree of job $j$ \\
$\pi_j,\ \delta_j$ & pickup and delivery operations of job $j$ &
  $\epsilon^{(l)},\ L$ & learnable GIN self-weight; number of GIN layers \\
$Q$ & per-type AMR carrying capacity &
  $\bm{e}_{\omega}$ & learned operation-type embedding \\
$\bm{a},\ a_j$ & AMR assignment vector; AMR assigned to job $j$ &
  $d(\omega,j)$ & dock served by the action ($o_j$ if pickup, else $d_j$) \\
$\sigma,\ \sigma_t$ & operation sequence; prefix committed up to step $t$ &
  $z_{\omega,k,j}$ & logit of action $(\omega,k,j)$ \\
$s^{\pi}_j,c^{\pi}_j$ & start / completion of pickup service of job $j$ &
  $V_{\phi}$ & critic (state-value) network, parameters $\phi$ \\
$s^{\delta}_j,c^{\delta}_j$ & start / completion of delivery service of job $j$ &
  & \\
$t_k(u,v)$ & realized travel time of AMR $k$ from $u$ to $v$ &
  \multicolumn{2}{@{}l}{\itshape Training} \\
$\Phi$ & collision-aware executor &
  $R,\ R(\sigma_{1:t})$ & dispatching rule; completion of prefix $\sigma_{1:t}$ by $R$ \\
$C_{\max},\ F_k$ & makespan; finish time of AMR $k$ incl.\ depot return &
  $C[\cdot]$ & shorthand for $C_{\max}(\Phi(\cdot))$ \\
$W_d$ & number of physical waiting cells at dock $d$ &
  $A_t,\ \hat{A}_t$ & stepwise counterfactual advantage; standardized \\
\multicolumn{2}{@{}l}{\itshape Sequential decision process} &
  $A^{\mathrm{lb}}_t$ & load-balance shaping advantage \\
$s_t,\ a_t$ & state and action at decision epoch $t$ &
  $\lambda_{\mathrm{lb}},\ \lambda_{H}$ & shaping and entropy coefficients \\
$\omega$ & operation type, $\omega\in\{\textsc{pick},\textsc{deliver}\}$ &
  $\mathcal{H}_t$ & entropy of the masked policy at step $t$ \\
$\Lambda(s_t)$ & feasible-action set (mask) in state $s_t$ &
  $\mathcal{L}(\theta)$ & training loss \\
$\hat{\Psi}$ & calibrated surrogate transition model &
  $B,\ E$ & batch size (episodes); number of training epochs \\
$\mathcal{E}_d$ & committed service events $(s_e,c_e,j_e,\omega_e)$ at dock $d$ &
  $T_{\mathrm{val}},\ \kappa$ & validation interval; infeasibility penalty \\
$\hat{b}_k,\ x_k$ & next-available time and position of AMR $k$ &
  $\mathbf{1}[\cdot]$ & indicator function \\
$\rho$ & ready time of an AMR at a dock &
  & \\
$\hat{t}(u,v)$ & calibrated travel-time estimate \eqref{eq:travelcal} &
  & \\
$\hat{w}_d(\rho)$ & queue-aware dock-wait estimate \eqref{eq:waitcal} &
  & \\
$q_d(\rho)$ & services still unfinished at time $\rho$ at dock $d$ &
  & \\
$\alpha,\ \beta$ & travel-calibration scale and offset &
  & \\
$\varphi(\cdot)$ & dock-wait penalty table over queue depths &
  & \\
$\pi_{\theta}$ & scheduling policy with parameters $\theta$ &
  & \\
\bottomrule
\end{tabular}
\end{table*}

\subsection{System Description}

We consider the internal transfer problem of a cross-docking terminal
served by a fleet of autonomous mobile robots (AMRs). Table~\ref{tab:notation}
summarizes the notation used throughout the paper. The terminal floor
is modeled as a four-connected unit grid $\mathcal{W}$ in which a set of
\emph{inbound docks} $\mathcal{D}^{\mathrm{in}}$ is located along one
boundary and a set of \emph{outbound docks} (processing stations)
$\mathcal{D}^{\mathrm{out}}$ along the opposite boundary; the AMR depots
lie between them. Consistent with cross-docking practice, the floor is
free of static obstacles, so congestion arises from vehicle interference
and dock contention rather than from the layout itself. AMRs travel at
unit speed between adjacent cells, and the distance between two
positions $u,v$ is the Manhattan metric $\Delta(u,v)$.

A set of transfer jobs $\mathcal{J}=\{1,\dots,n\}$ must be moved across
the terminal by a fleet $\mathcal{M}=\{1,\dots,m\}$ of homogeneous AMRs.
Each job $j\in\mathcal{J}$ represents one unit load of a material type
$\tau_j\in\mathcal{T}$ that becomes available at its inbound dock
$o_j\in\mathcal{D}^{\mathrm{in}}$ at release time $r_j\ge 0$ and must be
delivered to a designated outbound dock
$d_j\in\mathcal{D}^{\mathrm{out}}$. Handling a unit load of type $\tau$
occupies a dock for a type-dependent service duration $p(\tau)$; we
write $p_j=p(\tau_j)$ and assume the loading service at the inbound dock
and the unloading/processing service at the outbound dock take the same
time, reflecting symmetric handling of a unit load. Each job therefore
consists of an ordered pair of operations: a \emph{pickup} operation
$\pi_j$ executed at $o_j$ and a \emph{delivery} operation $\delta_j$
executed at $d_j$, both performed by the same AMR (no transshipment).
An AMR may consolidate loads subject to a per-type carrying capacity
$Q$: at no time may it carry more than $Q$ units of any single material
type.

Docks are exclusive single-server resources: at most one AMR may be in
service at a dock at any time, and services are non-preemptive. An AMR
arriving at an occupied dock queues in a bounded set of waiting cells
adjacent to that dock; when the waiting line is saturated, the AMR must
hold at an upstream position and re-approach later. Finally, AMR motion
must be collision-free: at any unit time step, each grid cell may be
occupied by at most one AMR.

\subsection{Formal Statement}

A solution is a pair $(\bm{a},\sigma)$, where
$\bm{a}=(a_1,\dots,a_n)\in\mathcal{M}^{n}$ assigns each job to an AMR
and $\sigma$ is a total order over the $2n$ operations
$\{\pi_j,\delta_j\}_{j\in\mathcal{J}}$ that determines the sequence in
which operations are committed. Let $s^{\pi}_j,c^{\pi}_j$
(resp.\ $s^{\delta}_j,c^{\delta}_j$) denote the service start and
completion times of $\pi_j$ (resp.\ $\delta_j$) induced by executing
$(\bm{a},\sigma)$. A feasible execution satisfies:
\begin{align}
c^{\pi}_j &= s^{\pi}_j + p_j, \qquad
c^{\delta}_j = s^{\delta}_j + p_j,
&& \forall j\in\mathcal{J}, \label{eq:nonpreempt}\\
s^{\pi}_j &\ge r_j,
&& \forall j\in\mathcal{J}, \label{eq:release}\\
s^{\delta}_j &\ge c^{\pi}_j,
&& \forall j\in\mathcal{J}, \label{eq:precedence}
\end{align}
\begin{align}
\big|\{\,j : a_j{=}k,\ \tau_j{=}\tau,\ c^{\pi}_j \le u < c^{\delta}_j \}\big| &\le Q,
&& \forall k,\tau,u, \label{eq:capacity}\\
\big[s_e, c_e\big) \cap \big[s_{e'}, c_{e'}\big) &= \emptyset,
&& \forall e\ne e' \ \text{at the same dock}, \label{eq:exclusive}
\end{align}
where \eqref{eq:nonpreempt} enforces non-preemptive services,
\eqref{eq:release} release dates, \eqref{eq:precedence} pickup-before-%
delivery precedence, \eqref{eq:capacity} the per-type carrying capacity,
and \eqref{eq:exclusive} dock exclusivity over all service events $e$.
In addition, for any two operations $u\prec v$ executed consecutively by
the same AMR $k$,
\begin{equation}
s_v \;\ge\; c_u + t_k\big(\mathrm{loc}(u),\mathrm{loc}(v)\big),
\label{eq:travel}
\end{equation}
where $t_k(\cdot,\cdot)\ge\Delta(\cdot,\cdot)$ is the realized travel
time of AMR $k$, which may exceed the Manhattan distance due to
collision avoidance, queuing detours, and dock-clearing maneuvers.

The realized travel and waiting times are not fixed constants but are
determined by a high-fidelity \emph{executor} $\Phi$, which routes every
AMR leg with space--time $\mathrm{A}^{*}$ search over a shared
reservation table (at most one AMR per cell per time step), manages the
bounded dock waiting lines, and inserts upstream holding maneuvers when
lines are saturated. The problem is therefore a simulation-based
optimization: find
\begin{equation}
(\bm{a}^{*},\sigma^{*}) \;=\;
\arg\min_{(\bm{a},\sigma)} \; C_{\max}\big(\Phi(\bm{a},\sigma)\big),
\label{eq:objective}
\end{equation}
where the makespan $C_{\max}=\max_{k\in\mathcal{M}} F_k$ and $F_k$ is
the time at which AMR $k$ completes its final activity, including its
return trip to the depot. The problem generalizes both the capacitated
pickup-and-delivery problem and parallel-machine scheduling with
sequence-dependent setups, and is therefore NP-hard; the coupling to
collision-free multi-agent routing further compounds its difficulty
\cite{stern2019mapf}.

\subsection{Sequential Decision Formulation}
\label{sec:mdp}

We cast the construction of $(\bm{a},\sigma)$ as an episodic sequential
decision process with exactly $2n$ decision epochs, one per committed
operation. At epoch $t$, the state
\begin{equation}
s_t=\big(\{\bm{x}^{A}_{k}\},\, \{\bm{x}^{J}_{j}\},\, \{\bm{x}^{D}_{d}\},\,
\{\mathcal{E}_d\},\, \sigma_{t}\big)
\label{eq:state}
\end{equation}
collects the AMR states (position, next-available time, per-type
inventory), the job states (unpicked / onboard / delivered, carrier
identity), the dock states, the set of committed dock service events
$\mathcal{E}_d=\{(s_e,c_e,j_e,\omega_e)\}$ at each dock $d$, and the
partial operation sequence $\sigma_t$.

An action is a triple
\begin{equation}
a_t=(\omega,k,j)\in\{\textsc{pick},\textsc{deliver}\}\times\mathcal{M}\times\mathcal{J},
\label{eq:action}
\end{equation}
giving a flat action space of size $2mn$ restricted by a feasibility
mask $\Lambda(s_t)$: $(\textsc{pick},k,j)$ is legal iff $j$ is unpicked
and AMR $k$ carries fewer than $Q$ units of type $\tau_j$;
$(\textsc{deliver},k,j)$ is legal iff $j$ is onboard AMR $k$. The mask
guarantees by construction that every complete episode yields a feasible
solution, so no repair mechanism is required.

Transitions follow a deterministic surrogate model $\hat{\Psi}$ that
advances the acting AMR's clock using a \emph{calibrated} travel-time
estimate $\hat{t}(\cdot,\cdot)$ and a queue-aware dock-wait estimate
$\hat{w}_d(\cdot)$ fitted offline against the executor $\Phi$
(Section~\ref{sec:calibration}); committing an action appends the
corresponding service window to $\mathcal{E}_d$. For a pickup, the
service start of job $j$ by AMR $k$ located at $x_k$ with next-available
time $\hat{b}_k$ is
\begin{equation}
s^{\pi}_j=\underbrace{\max\!\big(\hat{b}_k+\hat{t}(x_k,o_j),\, r_j\big)}_{\text{ready time }\rho}
\;+\;\hat{w}_{o_j}(\rho),
\label{eq:surrogate}
\end{equation}
and analogously for deliveries at $d_j$. After $2n$ steps the completed
solution is evaluated by the executor, and the learning objective is to
minimize the expected executed makespan
$\mathbb{E}_{\pi_\theta}\!\left[C_{\max}(\Phi(\bm{a},\sigma))\right]$;
the policy-gradient training procedure and its dispatch-rule baseline
are described in Section~\ref{sec:training}.

% =========================================================================

\section{Dock-Aware Graph Scheduler}
\label{sec:architecture}

\subsection{Overview and Design Rationale}

Existing learning-to-dispatch architectures for shop scheduling encode
jobs (operations) as graph nodes and machines as flat feature vectors
\cite{zhang2020l2d,song2023fjsp}. Applied to cross-docking, this design
has three blind spots. First, the binding resources are the docks: with
$m$ AMRs contending for $|\mathcal{D}^{\mathrm{in}}|+|\mathcal{D}^{\mathrm{out}}|$
single-server docks, queuing at docks---not processing---dominates the
makespan, yet job-only graphs represent dock congestion at best as
scalar features replicated across jobs. Second, the temporal structure
of dock occupancy (which service windows are already committed, and
when) is invisible. Third, the AMR fleet is encoded independently
per vehicle, so mutual interference cannot be anticipated.

The proposed scheduler, hereafter \textsc{DockGNN}, addresses all three
by treating docks as first-class graph entities. It comprises
(i)~heterogeneous encoders for AMR, job, and dock nodes;
(ii)~an AMR--dock interaction block that exchanges information between
the fleet and the dock set;
(iii)~a job \emph{contention--precedence} graph processed by
degree-normalized graph isomorphism network (GIN) layers
\cite{xu2019gin}; and
(iv)~a factored actor head that scores each masked action from the
operation type, the acting AMR, the target job, the operation-specific
dock, and a small set of per-action decision features. The full model
has $6.2\times10^{5}$ parameters ($5.7\times10^{5}$ actor,
$4.9\times10^{4}$ critic) with hidden width $h=128$ throughout.

% TODO: architecture overview figure
% \begin{figure*}[t]
%   \centering
%   \includegraphics[width=\textwidth]{figures/dockgnn_architecture}
%   \caption{Overview of \textsc{DockGNN}. Dock, AMR, and job features are
%   encoded separately; an AMR--dock attention block exchanges congestion
%   information; jobs fuse their inbound/outbound dock embeddings and pass
%   messages over the contention--precedence graph; the actor scores each
%   feasible (operation, AMR, job) triple.}
%   \label{fig:architecture}
% \end{figure*}

\subsection{Heterogeneous State Representation}
\label{sec:features}

At every decision epoch the state \eqref{eq:state} is rendered into four
feature families, summarized in Table~\ref{tab:features}. All features
are normalized to comparable magnitudes by fixed scales.

\emph{Dock nodes.} Each dock $d$ maintains its committed service events
$\mathcal{E}_d$. Writing $t$ for the current decision time, the derived
features are the availability delay
$\max(0,\mathrm{avail}_d-t)$, the number of queued future services
$q_d=|\{e\in\mathcal{E}_d : s_e>t\}|$, the waiting-line saturation
$\min(q_d/W_d,1)$ with $W_d$ the number of physical waiting cells, the
residual time of the service in progress, and the committed workload
$\sum_{e\,\text{unfinished}}\big(c_e-\max(s_e,t)\big)$, together with
the dock class (inbound/outbound) and its coordinates.

\emph{AMR nodes.} Each AMR encodes its busy flag and time-to-free,
per-type inventory ratios and residual capacity, the number of loads
currently onboard, spatial context (coordinates, distance to the
nearest dock, a near-dock indicator), and a local fleet density---the
fraction of other AMRs within a two-cell radius---which is the model's
proxy for imminent interference.

\emph{Job nodes.} Each job encodes its service duration, waiting time
since release, lifecycle status (unpicked/onboard/delivered), carrier
identity, material type, dock indices, and a completion-time estimate
computed with the calibrated travel model
(Section~\ref{sec:calibration}).

\emph{Per-action features.} For every candidate triple $(\omega,k,j)$,
three decision-grade quantities are computed directly: the calibrated
travel time from AMR $k$'s position to the service dock, the calibrated
downstream distance ($o_j$ to $d_j$ for pickups; zero for deliveries),
and the projected queue-aware dock wait $\hat{w}(\cdot)$ of
\eqref{eq:surrogate}. Providing these estimates explicitly relieves the
network from recovering travel arithmetic from coordinates through
several nonlinear layers.

\begin{table}[t]
\caption{Node and Action Features of \textsc{DockGNN}}
\label{tab:features}
\centering
\footnotesize
\begin{tabular}{@{}llc@{}}
\toprule
& Feature & Scale \\
\midrule
\multirow{9}{*}{\rotatebox{90}{AMR $\bm{x}^{A}_k\!\in\!\mathbb{R}^{12}$}}
& busy indicator & -- \\
& time until available & $1/100$ \\
& inventory ratio, per type ($\times3$) & $1/Q$ \\
& residual capacity ratio & $1/(|\mathcal{T}|Q)$ \\
& loads onboard & $1/10$ \\
& near-dock indicator & -- \\
& distance to nearest dock & $1/20$ \\
& local fleet density ($\le\!2$ cells) & $1/(m{-}1)$ \\
& position $(x,y)$ ($\times2$) & $1/10$ \\
\midrule
\multirow{9}{*}{\rotatebox{90}{Job $\bm{x}^{J}_j\!\in\!\mathbb{R}^{11}$}}
& existence flag & -- \\
& service duration $p_j$ & $1/25$ \\
& waiting time since release & $1/100$ \\
& carrier index (0 if unpicked) & $1/m$ \\
& lifecycle status $\{0,\tfrac12,1\}$ & -- \\
& material type one-hot ($\times3$) & -- \\
& inbound dock index & $1/|\mathcal{D}^{\mathrm{in}}|$ \\
& outbound dock index & $1/|\mathcal{D}^{\mathrm{out}}|$ \\
& completion-time estimate & $1/500$ \\
\midrule
\multirow{7}{*}{\rotatebox{90}{Dock $\bm{x}^{D}_d\!\in\!\mathbb{R}^{9}$}}
& inbound / outbound class ($\times2$) & -- \\
& availability delay & $1/100$ \\
& queued services $q_d$ & $1/m$ \\
& waiting-line saturation & -- \\
& residual service in progress & $1/100$ \\
& committed workload & $1/500$ \\
& position $(x,y)$ ($\times2$) & $1/10$ \\
\midrule
\multirow{3}{*}{\rotatebox{90}{Act.\ $\mathbb{R}^{3}$}}
& calibrated travel to service dock & $1/20$ \\
& calibrated downstream distance & $1/20$ \\
& projected queue-aware dock wait & $1/100$ \\
\bottomrule
\end{tabular}
\end{table}

\subsection{Entity Encoders and AMR--Dock Interaction}
\label{sec:interaction}

AMR and dock features pass through two-layer perceptrons and job
features through a linear map,
\begin{equation}
\bm{h}^{A}_k=\mathrm{MLP}_{A}(\bm{x}^{A}_k),\quad
\bm{h}^{D}_d=\mathrm{MLP}_{D}(\bm{x}^{D}_d),\quad
\bar{\bm{h}}^{J}_j=\bm{W}_{\!J}\,\bm{x}^{J}_j,
\end{equation}
all in $\mathbb{R}^{h}$. Fleet--dock coupling is then established by a
single pre-normalization residual attention round
\cite{vaswani2017attention,xiong2020preln} consisting of three
sublayers: docks attend to the fleet, the fleet attends to itself, and
the fleet attends to the updated docks,
\begin{align}
\bm{H}^{D} &\leftarrow \bm{H}^{D} +
\mathrm{MHA}\big(\mathrm{LN}(\bm{H}^{D}),\,\mathrm{LN}(\bm{H}^{A}),\,\mathrm{LN}(\bm{H}^{A})\big),
\label{eq:dock-from-amr}\\
\bm{H}^{A} &\leftarrow \bm{H}^{A} +
\mathrm{MHA}\big(\mathrm{LN}(\bm{H}^{A}),\,\mathrm{LN}(\bm{H}^{A}),\,\mathrm{LN}(\bm{H}^{A})\big),
\label{eq:amr-self}\\
\bm{H}^{A} &\leftarrow \bm{H}^{A} +
\mathrm{MHA}\big(\mathrm{LN}(\bm{H}^{A}),\,\mathrm{LN}(\bm{H}^{D}),\,\mathrm{LN}(\bm{H}^{D})\big),
\label{eq:amr-from-dock}
\end{align}
with four heads per attention. Through
\eqref{eq:dock-from-amr}--\eqref{eq:amr-from-dock}, a dock's embedding
reflects which vehicles are loaded and heading toward it \emph{before}
jobs consume that embedding, and each AMR perceives both the remaining
fleet and network-wide dock congestion.

\emph{Stability.} The output projections of all three attention
sublayers are zero-initialized, making the entire block an exact
identity at the start of training in the spirit of residual-branch
zero initialization \cite{bachlechner2021rezero,goyal2017accurate}.
This is not cosmetic: the embedding stream entering the block has small
magnitude relative to the unit-scale outputs that randomly initialized
attention produces from layer-normalized inputs, and without
zero initialization the three stacked residual additions amplify the
fleet embeddings by more than an order of magnitude at initialization,
which we found to prevent policy-gradient training from converging.

\subsection{Job Contention--Precedence Graph}
\label{sec:graph}

Jobs interact through two mechanisms with different lifetimes:
\emph{future contention} for shared docks, and \emph{committed
precedence} induced by scheduling decisions already made. Both are
encoded in a single dynamic graph $\mathcal{G}_t=(\mathcal{J},E_t)$
with symmetric edges over the job nodes:
(i)~a \emph{contention edge} joins every pair of unfinished jobs that
share an inbound or an outbound dock, so resource competition is
visible from the first decision onward; and
(ii)~\emph{order edges} join consecutively scheduled jobs on the same
AMR and on the same outbound dock, incrementally densifying the graph
as $\sigma_t$ grows.

Before message passing, each job fuses the embeddings of its two docks,
\begin{equation}
\bm{g}^{(0)}_j=\mathrm{MLP}_{F}\!\big(\big[\,\bar{\bm{h}}^{J}_j \,\|\,
\bm{h}^{D}_{o_j} \,\|\, \bm{h}^{D}_{d_j}\big]\big),
\label{eq:fusion}
\end{equation}
so every job representation is conditioned on the live congestion state
of exactly the two resources it will occupy. Message passing then
applies $L=3$ residual GIN layers with \emph{degree-normalized}
aggregation and layer normalization,
\begin{equation}
\bm{g}^{(l+1)}_j=\bm{g}^{(l)}_j+
\mathrm{LN}\!\left(\mathrm{MLP}^{(l)}\!\Big(
(1{+}\epsilon^{(l)})\,\bm{g}^{(l)}_j+
\tfrac{1}{\deg_j}\!\!\sum_{i\in\mathcal{N}(j)}\!\bm{g}^{(l)}_i
\Big)\right)\!.
\label{eq:gin}
\end{equation}
Mean aggregation deliberately trades the injectivity of sum-based GIN
\cite{xu2019gin} for scale invariance: contention edges form dock
cliques whose size grows linearly with $n/|\mathcal{D}|$, and
sum aggregation would make message magnitudes---and hence the
activation statistics of every downstream layer---depend on the
instance size, undermining transfer across job counts. With mean
aggregation the same trained network processes $n{=}10$ and $n{=}100$
instances at identical activation scales.

\subsection{Action Scoring and Value Estimation}
\label{sec:actor}

Let $\bm{e}_{\omega}\in\mathbb{R}^{h}$ be a learned embedding of the
operation type and let $d(\omega,j)=o_j$ if $\omega=\textsc{pick}$ and
$d_j$ otherwise denote the dock at which the candidate action would be
served. The actor scores every triple by
\begin{equation}
z_{\omega,k,j}=\mathrm{MLP}_{\pi}\!\big(\big[\,\bm{e}_{\omega}\,\|\,
\bm{h}^{A}_{k}\,\|\,\bm{g}^{(L)}_{j}\,\|\,
\bm{h}^{D}_{d(\omega,j)}\,\|\,\bm{x}^{\mathrm{act}}_{\omega,k,j}\big]\big),
\label{eq:logit}
\end{equation}
and the policy is the masked softmax
\begin{equation}
\pi_{\theta}(a\,|\,s_t)=
\frac{\exp(z_{a})\,\mathbf{1}[a\in\Lambda(s_t)]}
{\sum_{a'\in\Lambda(s_t)}\exp(z_{a'})}.
\label{eq:policy}
\end{equation}
The critic estimates the state value from pooled entity summaries,
\begin{equation}
V_{\phi}(s_t)=\mathrm{MLP}_{V}\!\big(\big[\,
\overline{\bm{g}}^{\mathrm{(unf)}} \,\|\, \overline{\bm{h}}^{A} \,\|\,
\overline{\bm{h}}^{D}\big]\big),
\label{eq:critic}
\end{equation}
where $\overline{\bm{g}}^{\mathrm{(unf)}}$ averages the embeddings of
unfinished jobs only. A full episode requires exactly $2n$ forward
passes; each pass scores $2mn$ actions in a single batched evaluation
of \eqref{eq:logit}, and all components are permutation-invariant in
the job and AMR sets, so the architecture applies unchanged to any
instance size.

% =========================================================================

\section{Surrogate Model Calibration}
\label{sec:calibration}

\subsection{The Two-Simulator Gap}

The policy constructs schedules on the surrogate $\hat{\Psi}$ of
Section~\ref{sec:mdp}, whereas solution quality is defined by the
executor $\Phi$ of \eqref{eq:objective}. An uncalibrated surrogate that
advances clocks by Manhattan distances and single-scalar dock
availabilities is systematically optimistic: it ignores collision
detours, dock-clearing maneuvers, waiting-line approach paths, and the
upstream holds triggered when waiting lines saturate. The state
features of Section~\ref{sec:features}---and therefore every decision
the policy makes---would then be conditioned on clocks that drift
further from reality as congestion grows, precisely in the regime where
scheduling decisions matter most. Rather than paying the computational
cost of embedding the full executor in the training loop, we close most
of this gap with two lightweight corrections fitted offline.

\subsection{Calibrated Travel and Dock-Wait Estimators}

\emph{Travel.} Realized leg times are modeled as an affine function of
the Manhattan distance, clamped from below by the metric itself:
\begin{equation}
\hat{t}(u,v)=\max\!\big(\Delta(u,v),\ \alpha\,\Delta(u,v)+\beta\big).
\label{eq:travelcal}
\end{equation}

\emph{Dock waiting.} The wait experienced by an AMR that becomes ready
at a dock $d$ at time $\rho$ is modeled as the committed availability
delay plus a penalty indexed by the queue depth at readiness,
\begin{equation}
\hat{w}_d(\rho)=\max\!\big(0,\ \mathrm{avail}_d-\rho\big)
+\varphi\!\big(\min(q_d(\rho),\,3)\big),
\label{eq:waitcal}
\end{equation}
where $q_d(\rho)=|\{e\in\mathcal{E}_d : c_e>\rho\}|$ counts committed
services unfinished at $\rho$, and
$\varphi:\{0,1,2,3^{+}\}\to\mathbb{R}_{\ge0}$ is a monotone lookup
table. The bounded table shape is deliberate: waiting-line detours and
upstream holds produce a cost cliff once the line saturates, a
nonlinearity that a linear congestion model cannot express.

\subsection{Fitting Procedure}

Calibration data are obtained by replaying complete schedules through
both simulators and aligning them per operation. Schedules are drawn
from heterogeneous sources---five dispatching-rule pairs and the
current learned policies---across instance sizes
$n\in\{10,20,40,60\}$, so that the fitted correction covers the state
distribution that policies actually visit. For every operation, the
raw surrogate supplies the Manhattan leg distance, the base wait, and
the queue depth at readiness; the executor's logged timeline is
partitioned, per AMR, into legs terminating at each service start,
within which movement entries sum to the realized travel and the
remainder is attributed to waiting. Executions containing infeasible
operations are discarded. The travel parameters $(\alpha,\beta)$ are
fitted by least squares and the penalties $\varphi(\cdot)$ as clipped
bin means with monotonicity enforced across depths.

In our environment, fitting on $10{,}360$ operation pairs from $159$
schedules yielded $\alpha=0.963$, $\beta=1.12$ ($R^{2}=0.91$) and
$\varphi=(0.04,\,0.04,\,1.66,\,3.57)$; the step between depths one and
two locates the waiting-line saturation cliff. The correction removes
the mean per-leg travel bias entirely ($+0.744\to0.000$ time units)
and halves the end-to-end makespan drift between surrogate and
executor (mean absolute error $24.3\to12.1$ on held-out schedules).
The calibration is a frozen artifact: it is fitted once before
training, adds no learned components to the training loop, and is
shared by \emph{all} policies and by the dispatching-rule baseline of
Section~\ref{sec:training}, so every method plans in the same
corrected world. After training converges, the fit can be re-estimated
from the new policies' schedules to confirm stability.

% =========================================================================

\section{Policy Optimization}
\label{sec:training}

\subsection{Objective and Gradient Estimator}

The policy $\pi_\theta$ is trained to minimize the expected executed
makespan $\mathbb{E}_{\pi_\theta}[C_{\max}(\Phi(\bm{a},\sigma))]$ with
the REINFORCE estimator \cite{williams1992reinforce}. During training,
episodes are generated by sampling from the masked policy
\eqref{eq:policy} on the calibrated surrogate; at validation and
deployment, actions are selected greedily. The gradient quality of
REINFORCE hinges on its baseline; we use a counterfactual baseline
built from a strong dispatching rule and evaluated by the executor
itself.

\subsection{Counterfactual Dispatch-Rule Baseline}
\label{sec:baseline}

Let $R$ denote a deterministic dispatching heuristic (job rule and AMR
rule; we use earliest-completion for both) and let
$R(\sigma_{1:t})$ denote the complete solution obtained by executing
the partial sequence $\sigma_{1:t}$ and completing all remaining
operations with $R$. Define the shorthand
$C[\cdot]=C_{\max}(\Phi(\cdot))$. The advantage of the action taken at
step $t$ is the improvement of committing that action and then
following the rule, relative to following the rule immediately:
\begin{equation}
A_t \;=\; C\big[R(\sigma_{1:t-1})\big]\;-\;
C\big[R(\sigma_{1:t-1}\oplus a_t)\big].
\label{eq:advantage}
\end{equation}
Both completions are evaluated by the collision-aware executor, so the
learning signal is grounded in true solution quality rather than in
surrogate clocks. Unlike episode-level rollout baselines
\cite{rennie2017self,kool2019attention}, \eqref{eq:advantage} assigns
credit per decision: an early pickup that congests a dock is penalized
at the step where it is committed, not diluted across the episode.
The price is $\mathcal{O}(n)$ executor evaluations per episode, which
dominates training cost; we consider the sharper credit assignment a
favorable trade at the instance sizes studied here, and
Section~\ref{sec:training-details} lists a cheaper actor--critic
alternative.

\subsection{Loss Function}

Advantages are standardized to zero mean and unit variance across all
steps of a batch of $B$ episodes. The per-batch loss combines the
policy gradient with a load-balancing shaping term and an entropy
bonus:
\begin{equation}
\mathcal{L}(\theta)=
-\frac{1}{B}\sum_{b=1}^{B}\sum_{t=1}^{2n}
\Big[\log\pi_\theta(a_t^b|s_t^b)\,
\big(\hat{A}_t^{\,b}+\lambda_{\mathrm{lb}}A_t^{\mathrm{lb},b}\big)
+\lambda_{H}\,\mathcal{H}_t^{b}\Big],
\label{eq:loss}
\end{equation}
where $\hat{A}_t$ is the standardized advantage,
$\mathcal{H}_t$ is the entropy of the masked policy over feasible
actions, and the shaping term
$A^{\mathrm{lb}}_t=\min_{k'}c_{k'}-c_{k}$ (for pickups; zero for
deliveries) compares the assignment count $c_k$ of the chosen AMR with
the least-loaded AMR, discouraging degenerate fleet utilization early
in training \cite{ng1999shaping}. We use
$\lambda_{\mathrm{lb}}=0.1$ and $\lambda_{H}=0.01$.

\begin{algorithm}[t]
\caption{Training with counterfactual dispatch baseline}
\label{alg:training}
\begin{algorithmic}[1]
\Require policy $\pi_\theta$, rule $R$, executor $\Phi$, surrogate $\hat{\Psi}$
\State fit calibration $(\alpha,\beta,\varphi)$ offline; freeze (Sec.~\ref{sec:calibration})
\For{$\mathrm{epoch}=1,\dots,E$}
  \For{$b=1,\dots,B$}
    \State sample instance; roll out $\sigma^b\!\sim\!\pi_\theta$ on $\hat{\Psi}$
    \For{$t=1,\dots,2n$}
      \State $A^b_t \gets C[R(\sigma^b_{1:t-1})]-C[R(\sigma^b_{1:t-1}\oplus a^b_t)]$
    \EndFor
  \EndFor
  \State standardize $\{A^b_t\}$ over the batch
  \State update $\theta$ by Adam on $\mathcal{L}(\theta)$ in \eqref{eq:loss}; clip $\lVert\nabla\rVert\le1$
  \If{$\mathrm{epoch}\bmod T_{\mathrm{val}}=0$}
    \State greedy-decode held-out instances; score by
    $C_{\max}+\kappa\cdot(\text{infeasible ops})$; keep best
  \EndIf
\EndFor
\end{algorithmic}
\end{algorithm}

\subsection{Training Protocol}
\label{sec:training-details}

Algorithm~\ref{alg:training} summarizes the procedure. Only the actor
parameters (encoders, interaction block, GIN layers, operation
embedding, and actor head) receive gradients under REINFORCE; the
critic \eqref{eq:critic} is reserved for an optional
proximal-policy-optimization variant \cite{schulman2017ppo} that
replaces \eqref{eq:advantage} with critic-based advantages when the
executor evaluations of the stepwise baseline are prohibitive. We
train with Adam \cite{kingma2015adam} (learning rate $3\times10^{-4}$,
batch $B=16$, gradient-norm clip $1.0$) for $2{,}000$ epochs. Every
$T_{\mathrm{val}}=50$ epochs the policy is evaluated by greedy decoding on a fixed
held-out validation set, disjoint from the training pool, and scored
by the executed makespan plus a penalty of $\kappa=10^{3}$ per
infeasible operation; the checkpoint with the best validation score is
retained. Two safeguards proved important in practice: the
identity-at-initialization residual branches of
Section~\ref{sec:interaction}, without which policy-gradient training
failed to converge, and computing the entropy in \eqref{eq:loss} over
the feasible-action support only, which keeps the bonus comparable as
the mask tightens toward the end of an episode.

% =========================================================================
% Suggested BibTeX entries (replace keys above with your bibliography):
%
% @inproceedings{vaswani2017attention, Vaswani et al., "Attention Is All
%   You Need", NeurIPS 2017}
% @inproceedings{xu2019gin, Xu, Hu, Leskovec, Jegelka, "How Powerful Are
%   Graph Neural Networks?", ICLR 2019}
% @inproceedings{xiong2020preln, Xiong et al., "On Layer Normalization in
%   the Transformer Architecture", ICML 2020}
% @inproceedings{bachlechner2021rezero, Bachlechner et al., "ReZero Is
%   All You Need: Fast Convergence at Large Depth", UAI 2021}
% @article{goyal2017accurate, Goyal et al., "Accurate, Large Minibatch
%   SGD: Training ImageNet in 1 Hour", arXiv:1706.02677}
% @inproceedings{zhang2020l2d, Zhang et al., "Learning to Dispatch for
%   Job Shop Scheduling via Deep Reinforcement Learning", NeurIPS 2020}
% @article{song2023fjsp, Song et al., "Flexible Job-Shop Scheduling via
%   Graph Neural Network and Deep Reinforcement Learning", IEEE TII 2023}
% @inproceedings{stern2019mapf, Stern et al., "Multi-Agent Pathfinding:
%   Definitions, Variants, and Benchmarks", SoCS 2019}
% @article{williams1992reinforce, Williams, "Simple Statistical
%   Gradient-Following Algorithms for Connectionist Reinforcement
%   Learning", Machine Learning 8, 1992}
% @inproceedings{rennie2017self, Rennie et al., "Self-Critical Sequence
%   Training for Image Captioning", CVPR 2017}
% @inproceedings{kool2019attention, Kool, van Hoof, Welling, "Attention,
%   Learn to Solve Routing Problems!", ICLR 2019}
% @inproceedings{ng1999shaping, Ng, Harada, Russell, "Policy Invariance
%   Under Reward Transformations: Theory and Application to Reward
%   Shaping", ICML 1999}
% @article{schulman2017ppo, Schulman et al., "Proximal Policy
%   Optimization Algorithms", arXiv:1707.06347}
% @inproceedings{kingma2015adam, Kingma, Ba, "Adam: A Method for
%   Stochastic Optimization", ICLR 2015}
% =========================================================================
